\documentclass[11pt,a4paper,twoside]{book}
\usepackage[T1]{fontenc}
\usepackage{lmodern}
\usepackage[utf8]{inputenc} 
\usepackage{amsmath, amsthm, amsfonts, amssymb, xfrac}
\usepackage{arydshln} %Linea dotted in array
\usepackage[italian]{babel}
\usepackage{titling} %Funzioni per modificare maketitle
\usepackage{titlesec} %Ridefinizione semplice di titoli 
\usepackage{float} %Comando H per posizione di float
\usepackage[left=2.5cm, right=2cm, top=2cm, bottom=2cm]{geometry} %Gestione margini pagina
\usepackage{perpage} %Permette di resettare una numerazione per pagina
\usepackage{calligra}
\usepackage{fancyhdr} %Lo uso per avere titolo e autore running
\usepackage[hidelinks]{hyperref} %lo uso per avere l'indice cliccabiile
\usepackage{graphicx, wrapfig} %Immagini
\usepackage[shortlabels]{enumitem} %Enumerate con diversi label, l'argomento mi permette di usare (a) al posto di label=(\alph*) etc
\usepackage{ulem}

\MakePerPage{footnote} %Resetta la numerazione delle footnote per pagina

%Autore-------------------------------------------
\author{Marco Accerenzi}
\date{}
\title{Appunti di fisica matematica}


%Titleformat-----------------------------------------------
\renewcommand{\thesection}{\arabic{section}}
\renewcommand{\thesubsection}{\thesection.\arabic{subsection}}
\titleformat{\section}{\huge\bfseries\boldmath}{{\Large \thesection}}{0.5em}{}
\titleformat{\subsection}{\Large\bfseries\boldmath}{{\large \thesubsection}}{0.5em}{}
\titleformat{\subsubsection}{\large\bfseries\boldmath}{{}}{0}{}
\titleformat{\paragraph}[runin]{\normalsize\bfseries}{{}}{0}{}

%Pagestyle pagine normali--------------------------------
\pagestyle{fancy} %Serve per poter modificare le pagine con fancyhdr facilmente
\renewcommand{\sectionmark}[1]{ \markright{#1}{} }

\lhead[]{\small Marco Accerenzi}  %Autore a sinistra nelle pagine dispari
\chead[\small \rightmark]{} %Capitolo al centro nelle pagine pari
\rhead[]{\small Appunti di Istituzioni di Metodi Matematici} %Titolo a destra nelle pagine dispari
\fancyfoot[LE,RO]{\large\thepage} %Numero pagina grande a sinistra in basso nelle pagine dispari
\fancyfoot[C]{}
\renewcommand{\headrulewidth}{0.25pt}
\renewcommand{\footrulewidth}{0.3pt}

%Modifica pagestyle pagine speciali-------------------------------------------
\fancypagestyle{plain}{%Ridefinisco il pagestyle della prima pagina dei chapter
\fancyhf{} %Toglie tutti i footer ed header
\fancyfoot[LO,RE]{\thepage} 
\renewcommand{\headrulewidth}{0pt}
\renewcommand{\footrulewidth}{0pt}
}

%Tcolorbox per i teoremi, esempi ed esercizi---------------------------
\usepackage[skins, breakable, theorems]{tcolorbox}

%Box per proposizione
\newtcbtheorem[number within=section]{Teorema}{Teorema}{
	enhanced,
	boxrule=0pt,frame hidden,
	borderline west={2.5pt}{0pt}{black},
	colback=yellow!1!white,
	sharp corners,
	coltitle=black,
	attach title to upper={\noindent},
	fonttitle=\bfseries\large,
}{thm}

%Box per proposizione
\newtcbtheorem[use counter from=Teorema]{Proposizione}{Proposizione}{
	enhanced,
	boxrule=0pt,frame hidden,
	borderline west={2.5pt}{0pt}{black},
	colback=yellow!1!white,
	sharp corners,
	coltitle=black,
	attach title to upper={\noindent},
	fonttitle=\bfseries\large,
}{thm}

%Box per Definizione
\newtcbtheorem[use counter from=Teorema]{Definizione}{Definizione}{
enhanced,
	boxrule=0pt,frame hidden,
	borderline west={2.5pt}{0pt}{black},
	colback=yellow!1!white,
	sharp corners,
	coltitle=black,
	attach title to upper={\noindent},
	fonttitle=\bfseries\large,
}{defi}

%Box per Lemma
\newtcbtheorem[use counter from=Teorema]{Lemma}{Lemma}{
enhanced,
	boxrule=0pt,frame hidden,
	borderline west={2.5pt}{0pt}{black},
	colback=yellow!1!white,
	sharp corners,
	coltitle=black,
	attach title to upper={\noindent},
	fonttitle=\bfseries\large,
}{lemm}

%Box per Corollario
\newtcbtheorem[use counter from=Teorema]{Corollario}{Corollario}{
enhanced,
	boxrule=0pt,frame hidden,
	borderline west={2.5pt}{0pt}{black},
	colback=yellow!1!white,
	sharp corners,
	coltitle=black,
	attach title to upper={\noindent},
	fonttitle=\bfseries\large,
}{Cor}



%Ambienti per proposizione, definizione, lemma e corollario
\newenvironment{teorema}[1][\empty]{\begin{Teorema}{
\ifx\empty#1\relax\else\mdseries\itshape#1\\\fi
}{}}{\end{Teorema}}
\newenvironment{proposizione}[1][\empty]{\begin{Proposizione}{
\ifx\empty#1\relax\else\mdseries\itshape#1\\\fi
}{}}{\end{Proposizione}}  
\newenvironment{definizione}[1][\empty]{\begin{Definizione}{
\ifx\empty#1\relax\else\mdseries\itshape#1\\\fi
}{}}{\end{Definizione}}
\newenvironment{lemma}[1][\empty]{\begin{Lemma}{
\ifx\empty#1\relax\else\mdseries\itshape#1\\\fi
}{}}{\end{Lemma}}
\newenvironment{corollario}[1][\empty]{\begin{Corollario}{
\ifx\empty#1\relax\else\mdseries\itshape#1\\\fi
}{}}{\end{Corollario}}

%Box per pezzi generici importanti-------------------------
\newenvironment{importante}[0]{ 
  \begin{center}\begin{tcolorbox}[enhanced,
  before skip=2mm,after skip=3mm,
  boxrule=0.4pt,left=5mm,right=2mm,top=1mm,bottom=1mm,
  colback=yellow!50,
  colframe=yellow!20!black,
  sharp corners,arc=3mm,
  underlay={
  	\path[fill=red!50!black,draw=none] (interior.south west) rectangle 			node[white]{\Huge\bfseries !} ([xshift=4mm]interior.north west);
  },
  drop fuzzy shadow]
  }{
 \normalsize \end{tcolorbox}\end{center}
}

%Box per esempi----------------------
\newenvironment{esempio}[1][Esempio]{ 
  \begin{center}\begin{tcolorbox}[
  title={#1:~},
  width=0.9\linewidth,
  flush right,
  enhanced jigsaw, 
  colback=white, 
  frame hidden, 
  sharpish corners, arc=3mm,
  borderline west = {1pt}{0pt}{gray},
  borderline south = {1pt}{0pt}{gray},
  coltitle=black, 
  fonttitle=\small\bfseries, 
  attach title to upper, 
  breakable,
  top=2.5pt, bottom=0.3pt]\small
  }{
 \normalsize \end{tcolorbox}\end{center}
}

%Box per osservazioni----------------------
\newenvironment{osservazione}[1][Osservazione]{ 
  \begin{center}\begin{tcolorbox}[
  title={#1:~},
  width=0.9\linewidth,
  flush right,
  enhanced jigsaw, 
  colback=white, 
  frame hidden, 
  sharpish corners, arc=3mm,
  borderline west = {1pt}{0pt}{gray},
  borderline south = {1pt}{0pt}{gray},
  coltitle=black, 
  fonttitle=\small\bfseries, 
  attach title to upper, 
  breakable,
  top=2.5pt, bottom=0.3pt]\small
  }{
 \normalsize \end{tcolorbox}\end{center}
}

%Ambiente proof con riga a sinistra---------------------------
\newenvironment{dimo}[1][\proofname]{
	\begin{tcolorbox}[
	breakable, 
	enhanced, 
	interior hidden, 
	frame hidden, 
	borderline west = {1pt}{0pt}{gray}, 
	top=0pt, bottom=0pt, left=0pt]
\begin{proof}[#1:~]
}
{
\end{proof}
\end{tcolorbox}
}

%Stili teorema per soluzione esercizi--------------------------------------
\newtheoremstyle{esss}{\topsep}{\topsep}{\normalfont}{0pt}{\bfseries}{}{ }{\thmname{#1}\thmnumber{ #2}\thmnote{ #3}}
\theoremstyle{esss}

\newtheorem{esercizio}{Esercizio}[section]

\newtheoremstyle{dotless}{\topsep}{\topsep}{\normalfont}{0pt}{\bfseries}{}{ }{\thmname{#1}\thmnumber{ #2}\thmnote{ #3}}
\theoremstyle{dotless}

\newtheorem*{ris}{}
\newenvironment{risoluzione}[1][]{  
  \begin{ris}[\bfseries Risoluzione #1]$ $\\
  }{
  \end{ris}
}

%Miste------------------------------------------------------------------
\numberwithin{equation}{section} %Numerazione equazione per sezione
\renewcommand{\thefootnote}{\fnsymbol{footnote}} 
\hypersetup{linktoc=all}
\newcommand{\sectionbreak}{\clearpage}

%Indice--------------------------------------------------------
\usepackage{tocloft}
\renewcommand{\cftpartfont}{\large\bfseries}
\renewcommand{\cftsecfont}{\bfseries}


%Inizio------------------------------------------------------------------------------

\begin{document}
\frontmatter


%Copertina----------------------------------------------------------------------
\begin{titlepage}
	\begin{center}

	\begin{center}
		\centering
		Anno accademico 2019-2020\par
		\vspace{.05\textheight}
	\end{center}	

		\vspace{.1\textheight}
		{\LARGE\scshape \textsc{Marco Accerenzi}\par}
		\vspace{.1\textheight}
		\textsc{}\par
		\vspace{.05\textheight}
		{\Huge \textsc{Metodi matematici}\par}
		\vspace{.05\textheight}
		{\large Riassunto degli appunti\par}
		\vfill
		{\centering\large
hhh
  }
  \vspace{2cm} 
\end{center}
\end{titlepage}
\thispagestyle{plain}
%----------------------------------------------------------------------------------

\setcounter{page}{1}
\tableofcontents\clearpage\thispagestyle{plain}

%Parte 1-----------------------------------------------------------------
\pagestyle{fancy}

\cleardoublepage
\mainmatter\setcounter{page}{1}
\addcontentsline{toc}{part}{Analisi Complessa}
\thispagestyle{plain}
	\begin{center}
		\vspace{.1\textheight}
		{\LARGE\scshape \textsc{Prima parte}\par}
		\vspace{.1\textheight}
		\textsc{}\par
		\vspace{.05\textheight}
		{\Huge \textsc{Analisi Complessa}\par}
		\vfill

\end{center}
\fancyfoot[C]{\small Analisi Complessa}
%------------------------------------------------------------------


\section{Funzioni complesse}
\pagestyle{fancy}
Come sappiamo i numeri complessi $z\in\mathbb{C}$ si possono esprimere come
\[
z=x+iy \qquad x,y\in\mathbb{R}
\]
in modo univoco, questo permette di identificare $\mathbb{C}$ con $\mathbb{R}^2$ tramite la biiezione $z\leftrightarrow \vec{x}=(x,y)$.
La topologia, cioè ``l`insieme degli aperti'', di $\mathbb{R}^2$ e di $\mathbb{C}$ sarà la stessa 
\[
I=\{|z|=\sqrt{x^2+y^2}<R\}\subseteq\mathbb{C} \leftrightarrow I=\{|\vec{x}|<R\}\subseteq\mathbb{R}^2
\]
Questo ci permetterà anche di identificare le funzioni complesse di variabile complessa $\mathbb{C}\rightarrow\mathbb{C}$ con campi vettoriali $\mathbb{R}^2\rightarrow\mathbb{R}^2$.

Sia $D\subseteq\mathbb{C}$ un aperto e $f:D\rightarrow\mathbb{C}$ una funzione $z\mapsto f(z)$ allora potremo usare l'identificazione $f\mapsto\vec{U}=(u,v):D\rightarrow\mathbb{R}^2$ per riscrivere $f$ come
\[
f(x,y)=u(x,y)+iv(x,y)
\]
Non faremo differenze nell'indicare l'insieme $D\subseteq\mathbb{C}$ e il suo corrispettivo in $\mathbb{R}^2$ (che sostanzialmente coincidono).

Ricordiamo che numero complesso $z$ ha anche una scrittura polare
\[
z=|z|e^{i\theta} \qquad 0\leq\theta<2\pi
\]
e che per $z=x+iy\in\mathbb{C}$ se ne definisce il complesso coniugato $\overline{z}=x-iy$.

Le parti reali ed immaginaria di un numero si possono scrivere come 
\[
x=\dfrac{z+\overline{z}}{2} \qquad y=\dfrac{z-\overline{z}}{2i}
\]

\begin{esempio}
Per $f(z)=\dfrac{z+\overline{z}}{2}$ vediamo che $f(z)=x$ e quindi $u_f(x,y)=x$ e $v_f(x,y)=0$, invece per $g(z)=z\overline{z}$ avremo $z\overline{z}=x^2+y^2\in\mathbb{R}$ e quindi $u_g(x,y)=x^2+y^2$ e $v_g(x,y)=0$.
\end{esempio}


\begin{definizione}
Una funzione complessa di variabile complessa $f:D\rightarrow\mathbb{C}$ si dice
\begin{enumerate}[i.]
\item Continua in $z\in D$ se vale
\begin{equation}\label{derivComplessa}
\lim_{h\rightarrow 0} f(z+h)=f(z) \qquad h\in\mathbb{C}
\end{equation}
\item Olomorfa (o analitica) in $z\in D$ se esiste finito il limite 
\begin{equation}\label{derivComplessa}
\lim_{h\rightarrow 0} \dfrac{f(z+h)-f(z)}{h}
\end{equation}
Chiameremo questo limite, se esiste, la derivata complessa di $f$ che indicheremo con $f'(z)$ o $\dfrac{df}{dz}$.\\
\item Antiolomorfa in $z\in D$ se esiste finito il limite
\[
\lim_{h\rightarrow 0} \dfrac{f(z+h)-f(z)}{\overline{h}}
\]
\end{enumerate}
$f$ si dice olomorfa (o antiolomorfa) in un aperto $E\subseteq D$ se è olomorfa (o antiolomorfa) in ogni punto $z\in E$.
\end{definizione}

Sottolineiamo che i limiti di questa definizione sono complessi, pensandoli nel piano reale sono cioè limiti bidimensionali che devono esistere in ogni direzione.

Dalla definizione è chiaro che una funzione $f(z)$ è olomorfa se e solo se la funzione $\overline{f(z)}$ è antiolomorfa.

La condizione di olomorfia \eqref{derivComplessa} equivale a chiedere che esista un $\alpha\in\mathbb{C}$ per cui
\[
f(z+h)-f(z)-\alpha h = o(|h|)
\]
dove $\alpha$ sarà la $f'$ calcolata in $z$.


\begin{importante}
Se $f$ è olomorfa in $z$ allora è anche continua in $z$.
\end{importante}

\begin{teorema}
Per $D\subseteq\mathbb{C}$ aperto una funzione $f:D\rightarrow\mathbb{C}$ è olomorfa in $z=x+iy\in D$ se e solo se il campo associato $\vec{f}:D\rightarrow\mathbb{R}^2$ è differenziabile in $\vec{x}=(x,y)$ e $\vec{f}=(u,v)$ rispetta le condizioni di Cauchy-Riemann
\[
\dfrac{\partial u}{\partial x}=\dfrac{\partial v}{\partial y} \qquad
\dfrac{\partial u}{\partial y}=-\dfrac{\partial v}{\partial x}
\]
\end{teorema}

\begin{dimo}
Il campo $\vec{f}$ si dice differenziabile in $\vec{x}$ se esiste una matrice $A$ tale che $\vec{f}(\vec{x}+\vec{h})-\vec{f}(\vec{x})-A\vec{h}=o(|\vec{h}|)$ e sappiamo che in tal caso la matrice $A$ è \[
A=\left(\begin{array}{cc} \partial_x u & \partial_y u \\ \partial_x v & \partial_y v\end{array}\right)
\] e $\vec{h}=\left( \begin{array}{c}h_1 \\ h_2\end{array}\right)$.
\begin{itemize}
\item Se $f$ è olomorfa in $z$ allora $\alpha=f'(z)=a+ib$ e con $h=h_1+ih_2$ si ha $\alpha h=ah_1-bh_2+i(bh_1+ah_2)\equiv \left( \begin{array}{cc}a & -b \\ b & a\end{array}\right) \vec{h}=A h$ con $a=\partial_x u$ e $b= \partial_x v$  

Quindi $f$ olomorfa $\rightarrow$ $\vec{f}=(u,v)$  differenziabile e $\partial_x u=\partial_y v$ e $\partial_y u=-\partial_x v$.
\item Se valgono le condizioni di Cauchy-Riemann e $\vec{f}$ è differenziabile allora
\[
|\vec{f}(\vec{x}+\vec{h})-\vec{f}(\vec{x})-A\vec{h}|=|f(z+h)-f(z)-\alpha h| 
\]
con $\alpha=\partial_x u + i\partial_x v=f'(z)$ \qedhere
\end{itemize}
\end{dimo}


 


Definiamo 
\[ \partial=\dfrac{\partial_x-i\partial_y}{2} \qquad \overline{\partial}=\dfrac{\partial_x+i\partial_y}{2}
\]
che esplicitamente  sono
\begin{gather*}
\partial f = \dfrac{\partial_x u - i \partial_y u  +i\partial_x v +\partial_y v }{2} \\
\overline{\partial} f = \dfrac{\partial_x u + i \partial_y u  +i\partial_x v -\partial_y v }{2}
\end{gather*}
dalla dimostrazione del teorema sappiamo che 
\begin{equation}\label{eqDerivataComplessaEsplicita}
\partial f =f'=\partial_x u+i\partial_x v
\end{equation}
e che le condizioni di Cauchy-Riemann equivalgono a $\overline{\partial}f=0$.

Possiamo interpretare questo come dire che le funzioni olomorfe dipendo solo da $z$ mentre le funzioni antiolomorfe solo da $\overline{z}$.

\begin{osservazione}Per $g=\dfrac{\partial f}{\partial z}$ con $f$ olomorfa vediamo che \[\dfrac{\partial g}{\partial \overline{z}}=\dfrac{\partial }{\partial \overline{z}}\dfrac{\partial f}{\partial z}=\dfrac{\partial }{\partial z}\dfrac{\partial f}{\partial \overline{z}}=0\] e quindi la derivata complessa di una funzione olomorfa è ancora olomorfa e una funzione olomorfa ammette quindi derivate di ogni ordine. Dimostreremo questa proprietà in modo rigoroso più avanti.
\end{osservazione}

Senza dimostrarle elenchiamo alcune regole per il calcolo delle derivate complesse:\begin{itemize}[label=$\cdot$]
\item Se $f$ e $g$ sono funzioni olomorfe allora $\partial(f+g)=f'+g'$ e $\partial (fg)=f'g+fg'$
\item Se $f:D_f\rightarrow E_f$ e $g:D\rightarrow E_g$ sono funzioni olomorfe tali che $f(D_f)\subseteq D_g$ allora $g\circ f$ è olomorfa con $\partial g(f(z))=g'(f(z))\dot f'(z)$
\item $\partial\left( \sum_{n=0}^N c_n z^n\right)  = \sum_n^N nc_n z^{n-1}$ con dominio di olomorfia $\mathbb{C}$.
\item Per $f$ olomorfa in $D$ e $f(z)\neq 0$ per $z\in\hat{D}\subseteq D$ la reciproca $F=\sfrac{1}{f}$ è olomorfa in $\hat{D}$ con derivata complessa $F'=-\dfrac{f'}{f^2}$
\item Se $f$ e $g$ sono funzioni olomorfe allora $\dfrac{f}{g}$ è olomorfa in $D_f\cap D_g\cap\{z:g(z)\neq0\}$ con derivata complessa $\partial \dfrac{f}{g}=\dfrac{gf'-fg'}{g^2}$
\item $f(z)=e^z$ ha parte reale $u=e^x\cos{y}$ e immaginaria $v=e^x\sin{y}$ e quindi è olomorfa con $f'(z)=e^z$. Analogamente $\partial \sin{z}=\cos{z}$ e $\partial cos{z}=-\sin{z}$
\item Siano $f$ e $g$ funzioni olomorfe in $D$ con $f(z_0)=g(z_0)=0$ ma con $g'(z_0)\neq0$ per $z_0\in D$ allora $\lim_{z\rightarrow z_0} \dfrac{f(z)}{g(z)}=\lim_{z\rightarrow z_0} \dfrac{f'(z)}{g'(z)}$
\end{itemize}

\subsection{Trasformazioni conformi}

\begin{definizione}
Una mappa $\mathbb{R}^2\rightarrow\mathbb{R}^2$ si dice una trasformazione conforme se preserva tutti gli angoli
\end{definizione}

Le funzioni olomorfe hanno un'interpretazione geometrica come mappe $\mathbb{R}^2\rightarrow\mathbb{R}^2$ che preservano gli angoli, chiamate trasformazioni conformi.


Un numero complesso $z$ individua una direzione nel piano $\mathbb{R}^2$ ovvero per $z=\rho e^{i\theta}$ l'angolo $\theta$ è l'angolo tra le ascisse e la direzione di $z$.

La tangente ad una curva $z(t)=x(t)+iy(t)$ in $z_0=z(t_0)$ è un vettore $v(t_0)=\left|\dfrac{dz}{dt}\right|_{t=t_0}$ 


\begin{proposizione}
Una funzione olomorfa rappresenta una trasformazione conforme.
\end{proposizione}
\begin{dimo}
Consideriamo l'angolo tra due curve $\gamma_1:t\mapsto z_1(t)$ e $\gamma_2:t\mapsto z_2(t)$ in un punto $z_1(0)=a$.
Le curve avranno vettori tangenti dati da
\[
V_1=\left.\dfrac{ dz_1}{dt}\right|_{t=0}
\]
Applicando una funzione olomorfa $f(z)$ i trasformati dei vettori tangenti saranno
\[
V_1^*=\left.\dfrac{d}{dt} f(z_1(t))\right|_{t=0}=f'(a)V_1 \qquad V_2^2=f'(a)V_2
\]
$f'(a)=\rho e^{i\theta}$ cioè entrambi i vettori tangenti sono ruotati dalla $f$ di uno stesso $\theta$ e l'angolo compreso è invariato. \qedhere
\end{dimo}

\begin{esempio}
\begin{itemize}[label=$\cdot$]
\item $f(z)=a+bz$ è una {\bfseries similitudine}
\item $f(z)=\dfrac{R^2}{\overline{z}}$ è ``un' inversione rispetto al cerchio di raggio $R$'', manda la famiglia di cerchi e rette in se stessa.
\end{itemize}
\end{esempio}

\begin{proposizione}
Una mappa $z\mapsto f(z)$ è conforme se e solo se $f$ è olomorfa o antiolomorfa.
\end{proposizione}

\begin{teorema}[Gourat]
Se $f:D\rightarrow\mathbb{C}$ è olomorfa in $z\in D$ allora $f'(z)$ è continua.
\end{teorema}
\begin{corollario}
$f$ è olomorfa in $D$ $\iff$ $\vec{f}$ è $C^1$ e soddisfa le CCR.
\end{corollario}

\section{Integrali curvilinei di funzioni complesse}

Diamo alcune definizioni già note sulle curve

\begin{definizione}
Una curva $[a,b]\rightarrow\mathbb{C}$ complessa $z(s)=x(s)+iy(s)$ si dice
\begin{itemize}
\item $C^1$ a tratti se $x(s)$ e $y(s)$ sono $C^1$.
\item semplice se $z(s_1)\neq z(s_2)$ per $ s_1\neq s_2$ esclusi gli estremi.
\item chiusa se $z(a)=z(b)$.
\end{itemize}
La curva inversa di $\gamma(s):=z(s)$ è $-\gamma:{-b,-a}\rightarrow\mathbb{C}$ con $z(s)=z^I(-s)$
\end{definizione}

Da qui in avanti considereremo soltanto curve $C^1$ a tratti e semplici.

L'integrale curvilineo di una funzione complessa di variabile complessa $f:D\rightarrow\mathbb{C}$ lungo una curva $\gamma\subset D$ è
\begin{align*}
\int_\gamma f(z)dz &= \int_a^b f(z(s))z'(s) ds= \int_a^b [u(z(s) +iv(z(s))] \cdot [x'(s)+y'(s)] ds =\\
&=\int_a^b \left[ u(z(s)) x'(s) -v(z(s))y'(s) \right] ds + i\int_a^b \left[u(z(s))y'(s)+v(z(s))x'(s) \right] ds
\end{align*}

e quindi definendo le due forme differenziali
\[
\phi_1=udx-vdy \qquad \phi_2=udy+vdx
\]
l'integrale è
\[
\int_\gamma f dx = \int_\gamma \phi_1 +i\int_\gamma \phi_2
\]

L'integrale è indipendente dalla parametrizzazione di $\gamma$, per $z(s)\rightarrow \tilde{z}(s)=z(\alpha(s))$ si avrà $z'(s)ds\rightarrow \tilde{z}'(s)ds=\dfrac{d}{ds} z(\alpha(s)) \,ds=\dfrac{dz}{d\alpha}\dfrac{d\alpha}{ds}ds=dz$

Ricordiamo inoltre che:\begin{itemize}[label=$\cdot$]
\item $\int_{-\gamma}f dz=-\int_\gamma f dz$
\item $\int_{\gamma_1+\gamma_2}fdz=\int_{\gamma_1}fdz+\int_{\gamma_2}fdz$
\item {\bfseries Disuguaglianza di Darboux} 
\[
\left| \int_\gamma f(z) dz \right|\leq \left| f(z)\right| |dz| \leq \max_{z\in\gamma}{|f|} \cdot L(\gamma)
\]
\end{itemize}

\subsection{Teorema di Cauchy}

Per una curva chiusa $\gamma\subset D$ semplice vale il \textit{teorema della curva di Jordan}
Il \textit{teorema della curva di Jordan} dice che $\mathbb{R}^2\setminus\gamma$ consiste di due insiemi aperti connessi: uno interno a $\gamma$, che chiameremo $S$ nel seguito, limitato  e uno esterno, che chiamiamo $E$, illimitato. Inoltre $\partial S =\gamma=\partial E$ e $S$ sarà contraibile.

L'orientazione canonica di una $\gamma$ chiusa semplice sarà quella antioraria.

\begin{definizione}
Una curva chiusa $\gamma\subset D$ si dice contraibile in $D$ (o nullomotopa in $D$) se può essere deformata con continuità in $D$ ad un punto di $D$.
\end{definizione}
Questo equivale a dire che ''l'interno'' di una curva contraibile in $D$ è tutto contenuto in $D$, senza buchi.


Un aperto $D$ si dice semplicemente connesso se è connesso e se ogni curva chiusa $\gamma\subset D$ è contraibile in $D$.
\par
Semplicemente connesso e contraibile sono la stessa cosa in $\mathbb{R}^2$ ma non sono proprietà diverse in dimensione maggiore.
\begin{definizione}
Una 1-forma $\phi=\phi^x dx + \phi^y dy$ si dice \begin{enumerate}[i.]\item chiusa se $\partial_x \phi^y=\partial_y\phi^x$
\item esatta se $\phi=d\Lambda$ per una qualche $\Lambda:\mathbb{R}^2\rightarrow\mathbb{R}$ differenziabile. 
\end{enumerate}
\end{definizione}
Gli integrali delle forme esatte lungo una curva per la definizione di curva esatta dipendono solo dagli estremi della curva e quindi l'integrale lungo una curva chiusa di una 1-forma esatta è sempre nullo.

Una 1-forma esatta e $C^1$ sarà anche chiusa per il teorema di Schwarz.
\begin{teorema}[Teorema di Cauchy]
Sia $f:D\rightarrow\mathbb{C}$ olomorfa e $\gamma$ chiusa contraibile in $D$. Allora
\[
\oint_\gamma f(z) dz =0
\]
\end{teorema}
\begin{dimo}
Per 
\[
\int_\gamma f dz = \int_\gamma \phi_1 + i\int_\gamma \phi_2 \qquad \phi_1 = u dx - vdy\quad\phi_2=udy+vdx
\]
le $\phi_{1,2}$ sono chiuse se e solo se $\vec{f}$ rispetta le CCR.
Quindi se $f$ è olomorfa $\vec{f}$ è $C^1$ e sono verificate le CCR. Inoltre 
$\gamma=\partial S$ con $S$ contraibile quindi $\exists \hat{S}\subset D$ contraibile con $\gamma\subset \hat{S}$.
Per il \textit{Lemma di Poincarè} ogni 1-forma $C^2$ chiusa in $D$ contraibile è anche esatta. 
Quindi $\phi_{1,2}$ sono esatte e l'integrale è nullo. \qedhere
\end{dimo}
\begin{corollario}
Se una $\gamma=\gamma_1-\gamma_2$ è contraibile in $D\subset\mathbb{C}$ e $f:D\rightarrow\mathbb{C}$ è olomorfa allora
\[
\int_{\gamma_1} f(z) dz = \int_{\gamma_2} f(z) dz
\]
\end{corollario}
Quindi per calcolare l'integrale lungo una curva di una funzione olomorfa potremo deformare a piacere il cammino di integrazione nel dominio di olomorfia con sicurezza di ottenere lo stesso risultato.

\begin{corollario}
Sia $f$ olomorfa in $D$ e siano $\gamma_{1,2}\subset D$ due curve chiuse in $D$ tali che $\gamma_1=\partial S_1$ e $\gamma_2=\partial S_2$ con $S_2\subset S_1$\footnotemark. Allora se $S_1\setminus S_2\subset D$ 
\[
\oint_{\gamma_1} f(z)dz = \oint_{\gamma_2} f(z) dz
\] 
\end{corollario}
\footnotetext{Questo significa che potrebbero esserci punti di non olomorfia all'interno di $S_2$ e che quindi non si può applicare il teorema di Cauchy}
\begin{dimo}
$f$ è olomorfa in $A=S_1\setminus S_2$ che che contraibile e quindi $\partial A=\gamma_1+\Gamma-\gamma_2-\Gamma:=\gamma$ è una curva chiusa contraibili in $D$. Da questo segue che $\oint_\gamma f(z)dz=0=\oint_{\gamma_1}f\,dz-\oint_{\gamma_2}f\,dz=0$ \qedhere
\end{dimo}

\begin{esempio}
Per $f(z)=(z-a)^n$ con $n\in\mathbb{Z}$ $f$ è olomorfa in $\mathbb{C}\setminus\{a\}$ per ogni $n$. Calcoliamo $\oint_{\gamma}f(z)\,dz$. \begin{enumerate}[i.]
\item Se $\gamma$ non circonda $a$ dal teorema di Cauchy $\oint_{\gamma}f(z)\,dz=0$
\item Se $\gamma$ circonda $a$ (una sola volta) possiamo considerare la curva equivalente in $\mathbb{C}\setminus\{a\}$ data da $C_R(a)$ in cui $z=a+Re^{i\theta}$. Allora
\[
\oint_{C_R(a)}(z-a)^n \, dz= \oint_0^{2\pi} R^n e^{in\theta} iRe^{i\theta} d\theta=\begin{cases} \int_0^{2\pi} R^{n+1}i e^{i\theta(n+1)} d\theta=0 & n\neq -1 \\ \int_0^{2\pi} i d\theta =2\pi i & n=-1 \end{cases}
\]
\end{enumerate}
\end{esempio}

\subsection{Formula fondamentale del calcolo integrale}

\begin{teorema}[Teorema fondamentale del calcolo]
Data una funzione $f$ olomorfa in $D$ contraibile e uno $z_0\in D$ esiste un'unica funzione olomorfa $F$ tale che $F'(z)=f(z)$ e questa è
\[
F(z)-F(z_0)=\int_{\underset{\gamma}{z_0}}^z f(w) \, dw
\]
con $\gamma$ una curva arbitraria in $D$.
\end{teorema}

\begin{dimo}
$f(z) \,dz=\phi_1+i\phi_2$ dove $\begin{cases} \phi_1=udx-vdy \\ \phi_2=udy-vdx\end{cases}$ sono forme di classe $C^1$. 
L'olomorfia di $f$ implica che $\phi_1$ e $\phi_2$ siano forme chiuse in $D$ contraibile e quindi forme esatte
\[
\exists U,V:D\rightarrow\mathbb{C} : \phi_1=dU \quad \phi_2=dV
\]
Sia ora $F=U+iV$. Allora dato $f$ è contianua e che $\partial_x U=u=\partial_y V$ e $ \partial_x V=v=-\partial_y U$ $F$ è olomorfa in $D$.

Dalla \eqref{eqDerivataComplessaEsplicita} troviamo $F'=\partial_x F= u+iv=f$ e infine 

\[
\int_{z_0,\gamma}^z f\, dz=\int_a^b F'(z(s)) \dfrac{dz}{ds} ds = \int_a^b \dfrac{d}{ds} F(z(s))ds=F(z)-F(z_0) \qedhere
\]

\end{dimo}




\subsection{Formula di Cauchy}

Chiameremo $L^p(\mathbb{R}^D$ l'insieme $\{f:\mathbb{R}^D\rightarrow\mathbb{C}:\int_{\mathbb{R}^D} |f(\vec{x})|^p d^D x<\infty\}$.

Ricordiamo il teorema della convergenza dominata: per una successione $f_n(\vec{x})\in L^1$ $\forall n$ tale che esista il limite puntuale $\lim_{n\rightarrow\infty}f_n(\vec{x})=f(\vec{x})$ quasi ovunque e che esista una $g(\vec{x})\in L^1$ tale che $|f(\vec{x})|\leq g(\vec{x})$ $\forall n$ allora $f(\vec{x})\in L^1$ e ``posso portare il limite nell'intgrale'' cioè \[ \lim_{n}\int f_{n}d^D x=\int \lim_n f_n d^Dx = \int f d^D x\]  



\begin{teorema}[Formula di Cauchy]
Sia $f:D\rightarrow\mathbb{C}$ olomorfa in $D$ e $\gamma\subset D$ contraibile in $D$ e chiusa con $\gamma=\partial S$. Allora
\begin{equation}
\dfrac{1}{2\pi i}\oint_\gamma \dfrac{f(w)}{w-z}\, dw = \begin{cases} f(z) & \text{se } z\in S \\ 0 & z\not\in S \end{cases}
\end{equation}
\end{teorema}
La $\dfrac{f(w)}{w-z}$ è olomorfa in $D\setminus\{z\}$ quindi se $z\in S$ $\gamma$ non è contraibile nel suo dominio di olomorfia.
\begin{dimo}
Ovviamente se $z\not\in S$ la funzione integranda è olomorfa in $S$, $\gamma$ è contraibile nel dominio di olomorfia e per il teorema di Cauchy l'integrale è nullo.


Sia $z\in S$. $\gamma$ è equivalente in $D$ ad un circonferenza centrata in $z$ di raggio $R$ sufficientemente piccolo da essere contenuta in $S$.
\begin{align*}
\oint_\gamma \dfrac{f(w)}{w-z} dw&=\oint_{C_R} \dfrac{f(w)}{w-z} dw=\oint_{C_R} \dfrac{f(w)-f(z)}{w-z} dw+\oint_{C_R} \dfrac{f(z)}{w-z} dw=\\
&= \int_{0}^{2\pi} \dfrac{f(w)-f(z)}{Re^{i\theta}} Rie^{i\theta}\,d\theta+f(z)2\pi i=\int_{0}^{2\pi} \left(f(w)-f(z)\right) i\,d\theta + f(z)2\pi i
\end{align*}
Queste uguaglianze sono indipendenti da $R$ a patto che sia abbastanza piccolo quindi posso anche prenderle al limite $R\rightarrow0$. Vediamo che
\[
\left| \oint_{C_R} \dfrac{f(w)}{w-z} dw \right| = \int_{0}^{2\pi} \left|f(w)-f(z)\right| \,d\theta \leq  \sup_{w\in C_R}\left| f(w)-f(z) \right| 2\pi \underset{R\rightarrow 0}{\longrightarrow} 0
\]
quindi 
\[
\oint_\gamma \dfrac{f(w)}{w-z} dw = 2\pi i f(z) \qedhere
\]
\end{dimo}

\begin{corollario}
Una funzione olomorfa in $D\subset\mathbb{C}$ ammette derivate di ogni ordine in $D$. Se $\gamma$ è chiusa e contraibile in $D$ contente $z$ allora
\[
f^{(n)}(z)=\dfrac{n!}{2\pi i} \oint_\gamma \dfrac{f(w)}{(w-z)^{n+1}} \, dw
\]
\end{corollario}
\begin{dimo}

\end{dimo}

\begin{teorema}[Teorema di Morera]
Per $f:D\rightarrow\mathbb{C}$ continua con $\oint_\gamma f(z) dz =0 $ $ \forall \gamma\subset D$ contraibile in $D$ . Allora $f$ è olomorfa in $D$.
\end{teorema}
\begin{dimo}
\end{dimo}

\section{Funzioni armoniche}
\begin{definizione}
Una funzione $\phi:\mathbb{R}^2\rightarrow\mathbb{R}$ di classe $C^1$ si dice armonica se soddisfa
\[
\nabla^2\phi=\left( \dfrac{\partial^2}{\partial x^2} + \dfrac{\partial^2}{\partial y^2} \right) \phi =0
\]
\end{definizione}
Ancora usando la biiezione tra $\mathbb{R}^2$ e $\mathbb{C}$ vediamo che una funzione armonica così definita equivale ad una funzione armonica complessa.
\begin{proposizione}
La più generale funzione armonica complessa è somma di una funzione olomorfa e di una funzione antiolomorfa
\[
\phi=F(z)+G(\overline{z})
\]
\end{proposizione}
\begin{dimo}
\end{dimo}

Da questo segue che le funzioni armoniche sono $C^\infty$ e che per $f=u+iv$ olomorfa $u$ e $v$ sono sempre armoniche.

\section{Teorema di Liouville}
\begin{teorema}[Teorema della media]
Sia $f:D\rightarrow\mathbb{C}$ olomorfa e $B\subset D$ un disco centrato in $z$ di raggio $R$ con $\gamma=\partial B$. Allora $f(z)$ è la media di $f$ sulla $\gamma$ cioè
\[
f(z) = \dfrac{1}{2\pi}\int_0^{2\pi} f(z+Re^{i\theta}) d\theta 
\]
\end{teorema}
\begin{dimo}
Dalla formula di Cauchy \[ f(z)=\dfrac{1}{2\pi i} \oint_{\gamma} \dfrac{f(w)}{w-z}\,dw=\dfrac{1}{2\pi i} \int_0^{2\pi} \dfrac{f(z+Re^{i\theta})}{Re^{i\theta}}Rie^{i\theta} \,d\theta \qedhere \]
\end{dimo}

\begin{teorema}[Principio del massimo modulo]
Sia $f:D\rightarrow\mathbb{C}$ olomorfa su $D$ aperto. Se $f$ non è costante in $D$ allora non esiste nessun massimo relativo di $|f|$ in $D$ ovvero nessun $z_m\in D$ tc $|f(z)|\leq|f(z_m)|$ $\forall z in I_{z_0}$. 
\end{teorema}
\begin{dimo}
Sia per assurdo $|f(z)|\leq|f(z_0)|$ $\forall |z-z_0|<R$. Sia poi $r<R$ e $f(z_0)=\dfrac{1}{2\pi} \int_0^{2\pi} f(z_0+re^{i\theta}) s\theta$.
Allora \begin{align*}
|f(z_0)|\leq\int_0^{2\pi} |f(z_0+re^{i\theta})|d\theta\leq|f(z_0)|&\implies |f(z_0)|=\int_0^{2\pi} |f(z_0+re^{i\theta})|d\theta \\ &\implies \int_0^{2\pi} \left[ |f(z_0)|-|f(z_0+re^{i\theta})| \right] \, d\theta=0 \quad \forall r<R
\end{align*}
Ma quindi $|f(z)|=|f(z_0)|$ $\forall |z-z_0|<R$ che è assurdo, quindi $f$ non ha massimo assoluto in $D$.

Se la funzione non può avere massimo assoluto in un qualunque aperto in cui è olomorfa non potrà avere neanche massimi relativi dato che un massimo relativo è assoluto in un intorno aperto sufficientemente piccolo. \qedhere
\end{dimo}

\begin{corollario}
Sia $f$ è olomorfa in $D$ e continua in $\overline{D}$. Allora se $f$ non è costante in $D$ il massimo di $|f(z)|$ è sul bordo e 
\[
|f(z)|<M \quad \forall z\in  D
\]
\end{corollario}

Se $f(z)\neq 0$ $\forall z\in D$ la sua reciproca sarà olomorfa in $D$ e potremo applicare ad essa il principio del massimo modulo, quindi se $f$ non si annulla anche i suoi minimi assoluti saranno sulla frontiera.

\begin{definizione} Una $f:\mathbb{C}\rightarrow\mathbb{C}$ si dice intera se è olomorfa su tutto $\mathbb{C}$
\end{definizione}

\begin{teorema}[Teorema di Liouville]
Data $f(z)$ una funzione intera che, per $a\geq0$ $b> 0$ e $n\in\mathbb{N}$ , sia $f(z)\leq a|z|^n+b$ $\forall z$. Allora $f(z)$ è un polinomio di grado al più $n$.
\end{teorema}
\begin{dimo}
Con la formula per le derivate dal corollario alla formula di Cauchy
\[
f^{(n+1)}(z)=\dfrac{(n+1)!}{2\pi i} \oint_{C_R} \dfrac{f(w)}{(w-z)^{n+2}} dw
\]
dove $C_R$ è un cerchio di raggio $R$ qualsiasi (perchè $f$ è intera) e quindi $w=z+Re^{i\theta}$ e $w-z=Re^{i\theta}$. Quindi
\begin{align*}
0 \leq \left| f^{(n+1)}(z)\right| & < \left| \dfrac{(n+1)!}{2\pi i} \right| \oint_{C_R} \left|\dfrac{f(w)}{(w-z)^{n+2}} \right| dw = \dfrac{(n+1)!}{2\pi i}\int_0^{2\pi} \dfrac{\left|f(w)\right|}{\left|Re^{i\theta}\right|^{n+2}} \left| iRe^{i\theta} \right| d\theta = \\
& \leq \dfrac{(n+1)!}{2\pi i}\int_0^{2\pi} \dfrac{\left|f(w)\right|}{R^{n+1}} d\theta \leq  
\dfrac{(n+1)!}{2R\pi i}\int_0^{2\pi} \dfrac{a \left|w \right|^n +b}{R^{n}} d\theta \leq \\ & \leq \dfrac{(n+1)!}{2R\pi i}\int_0^{2\pi} \dfrac{a \left|z+ R \right|^n +b}{R^{n}} d\theta  \underset{R\rightarrow \infty}{\longrightarrow} 0
\end{align*}
e $|f^{(n+1)}(z)|=0\implies f^{(n+1)}(z)=0$ e quini $f(z)$ è un polinomio di grado $\leq n$. \qedhere
\end{dimo}

In particolare se si trova $f(z)\leq a$ $\forall z$ allora $f(z)$ sarà una costante in $\mathbb{C}$.

In realtà nella dimostrazione di questo teorema usiamo $f(z)<a|z|^n+b$ per $|z|>R$ con $R$ fissato, per funzioni intere però richiedere che la disuguaglianza valga in un intorno di infinito o in tutto $\mathbb{C}$ è uguale per il principio del massimo.

\begin{teorema}[Teorema fondamentale dell'algebra]
Un polinomio $P_n(z)$ di ordine $n$ ha almeno $1$ zero $a\in\mathbb{C}$ 
\[
P_N(z)=(z-a)G_{N-1}(z) \qquad G_{N-1}(z)\text{ è un polinomio di ordine } N-1
\]
\end{teorema}
\begin{dimo}
Sia per assurdo $P_N(z)\neq 0$ in ogni $z\in\mathbb{C}$. Allora la sua reciproca $f(z)=\dfrac{1}{P_N(z)}$ è intera.
Inoltre chiaramente $f(z)\underset{|z|\rightarrow \infty}{\longrightarrow} 0$ quindi $|f(z)|<M$ in tutto un intorno $|z|>R$ per un qualche $R$ fissato. Ma allora $f(z)$ dovrebbe essere costante per il teorema di Liouville, che è assurdo.
Quindi $P_N(z)$ ha almeno uno zero.
\end{dimo}
\begin{corollario}
Ogni polinomio $P_N(z)$ di ordine $N$ ha $N$ radici complesse 
\[
P_N(z)=k (z-a_1) \dots (z-a_N)
\]
\end{corollario}
\begin{dimo}
$P_N=c_N z^N +\dots + c_1 z + c_0$ ha almeno uno zero $a_1\in\mathbb{C}$ e $z=(z-a_1)+a_1$ quindi
\begin{align*}
P_N(z) &=c_N [(z-a_1)+a_1]^N +c_{N-1} [(z-a_1)+a_1]^{N-1}+\dots + c_1 [(z-a_1)+a_1] + c_0 \\ &= C_N (z-a_1)^N + \tilde{c}_{N_1} (z-a_1)^{N-1} + \dots \tilde{c}_0
\end{align*}
con $P_N(a_1)=0 \implies \tilde{c}_0=0$. Possiamo raccogliere  $C_N(z-a_1)$ ottenendo
\begin{align*}
P_N(z) &= C_N (z-a_1) \left(d_n (z-a_1)^N + d_{N-1} (z-a_1)^{N-1} + \dots d_1 \right) = \\ & = C_N (z-a_1) \left(\tilde{d}_{N-1} z^{N-1} + \tilde{d}
_{N-2} z^{N-2} + \dots \tilde{d}_0 \right)= c_N (z-a_1) G_{N-1}(z)
\end{align*}
dove $G_{N-1}$ è un polinomio di grado $N-1$ che avrà uno zero in $a_2$. Posso quindi ripetere il procedimento altre $N-1$ trovando
\[
P_N(z)=k (z-a_1)(z-a_2)\dots (z-a_N) \qedhere
\]
\end{dimo}
\section{Rappresentazione per serie}

\begin{definizione}
Sia $\{f_n(z)\}$ una successione di funzioni $D\rightarrow\mathbb{C}$\begin{itemize}[label=$\cdot$]
\item $f_n$ converge puntualmente ad una $f:D\rightarrow\mathbb{C}$ su $D$ se $\forall z\in D$ $\lim_{n\rightarrow\infty}|f_n(z)-f(z)|=0$
\item $f_n$ converge uniformamente ad una $f:D\rightarrow\mathbb{C}$ su $D$ se 
\[
\forall \epsilon>0 \> \exists N : \sup_{z\in D} |f_n(z)-f(z)| <\epsilon \quad \forall n>N
\]
\item Si dice che $f_n$ è una successione uniformemente di Cauchy in $D$ se
\[
\forall \epsilon>0 \>\exists N : \sup_{z\in D}|f_n(z)-f_m(z)|<\epsilon \quad \forall n,m>N
\]
\end{itemize}
\end{definizione}

Ricordiamo dei risultati già noti in generale nel contesto delle funzioni complesse di variabile complessa.

\begin{teorema} Sia $f_n:D\rightarrow\mathbb{C}$ una successione di funzioni con $D\subseteq\mathbb{C}$\begin{itemize}[label=$\cdot$]
\item  $f_n$ è uniformemente di Cauchy in $D$ se e solo se è uniformemente convergente in $D$.
\item Se le $f_n$ sono continue uniformemente convergenti in $D$ a $f$. Allora $f$ è continua in $D$.
\item Se le $f_n$ convergono uniformemente a una $f$ in $A$ e $\int_A f_n<\infty$ allora $lim_n \int_A f_n = \int_A \lim_n f_n =\int_A f$. In particolare questo vale anche per $A:=\gamma$ una curva.
\end{itemize}
\end{teorema}

\begin{definizione}Data una successione di funzioni $f_n$ consideriamo la serie $\sum_{n}f_n$ :
\begin{itemize}[i.]
\item Si dice che la serie $\sum_{n=0}^\infty f_n(z)$ converge ad una funzione $F(z)$, puntualmente o uniformemente, se la successione delle $s_n(z)=\sum_{m=0}^n f_m(z)$ converge a $F(z)$, rispettivamente puntualmente o uniformemente.
\item Si dice che la serie $\sum_{n=0}^\infty f_n(z)$ è assolutamente convergente, puntualmente o uniformemente, se $\sum_{n=0}^\infty |f_n(z)|$ converge, puntualmente o uniformemente. In tal caso anche $\sum_n f_n$ è convergente.
\end{itemize}
\end{definizione}

\begin{teorema}[Teorema di Weierstrass]
Sia $f_n:A\rightarrow\mathbb{C}$ una successione di funzioni olomorfe in $A$ compatto e continue su $A$ e $\partial A=\gamma$ dove $\gamma$ è $C^1$ a tratti. Sia inoltre $\sum_n f_n$ convergente uniformemente su $R$. Allora:\begin{enumerate}[i.]
\item $\sum_n f_n$ converge uniformemte su $A$
\item $f=\sum_n^\infty f_n$ è olomorfa all'interno di $A$
\item $f^{(k)}(z)=\sum_n^\infty f_n^{(k)} (z)$ $\forall k\geq 0$
\end{enumerate}
\end{teorema}

\begin{dimo}~\begin{enumerate}[i.] 
\item Per $n,m>N$ dato che $\sum_n f_n$ è di Cauchy si ha \[|s_n(z)-s_m(z)|_{z\in A}=\left|\sum_{p=m}^n f_p(z)\right|_{z\in A}\overset{\text{max. modulo}}{\leq} \left|\sum_{p=m}^n f_p(z)\right|_{z_0\in \gamma}<\epsilon\]

Quindi $s_n(z)$ è uniformemente di Cauchy in $A$ e dunque è uniformemente convergente in $A$ e $f$ è continua.

\item Sia $\Gamma$ contraibile all'interno di $A$, allora $\oint_\Gamma f\, dz=\oint_\Gamma\sum_n^\infty f_n \, dz=\sum_m^\infty \oint_{\Gamma} f_n\,dz=0$ perchè le $f_n$ sono olomorfe, quindi per il teorema di Morera $f$ è olomorfa in $A$.

\item Sia $z$  un punto interno ad $A$ e $\Gamma$ un cerchio centrato in $z$ di raggio $R$ sufficentemente piccolo da essere contenuto in $A$. Allora
\[g_n(w):=\dfrac{f_n(w)}{(w-z)^k}\overunderset{\text{unif.}}{w\in\Gamma}{\longrightarrow}\dfrac{f(w)}{(w-z)^k}
\]
quindi dall'espressione per le derivate dalla formula di Cauchy troviamo
\[
f^{(k)}(z)=\dfrac{k!}{2\pi i} \oint_{\Gamma}\dfrac{f(w)}{(w-z)^{k+1}}\, dw =\overset{\text{conv. unif.}}{=}\sum_{n=0}^\infty \dfrac{k!}{2\pi i} \oint_{\Gamma} \dfrac{f_n(w)}{(w-z)^{k+1}}\,dw=\sum_{n=0}^\infty f_n^{(k)}(z) \quad \text{puntualmente} \qedhere
\]
\end{enumerate}
\end{dimo}

\subsection{Serie di potenze}

Ricordiamo il già noto criterio di Weierstrass: se in un disco si ha $\sup_{|w|<\rho} |c_n w^n|\leq a_n$ con $\sum_n^\infty a_n<\infty$ allora la serie $\sum_n^\infty c_n w^n$ converge uniformemente nel disco.

Ovviamente se una serie $\sum_n c_n w^n$ diverge per un $w=b$ allora divergerà anche per ogni $w>b$.

\begin{teorema}[Teorema di Abel]
Se la serie di potenze $\sum_{n=0}^\infty c_n(z-a)^n$ converge per un $z_0\neq a$ allora converge assolutamente uniformemente su ogni disco chiuso di raggio $|z-a|\leq \rho <|z_0-a|$.
\end{teorema}


\begin{dimo}
Poniamo $w=z-a$ e  $w_0=z_0-a$, vogliamo mostrare che se $\sum_n^\infty c_nw_0^n$ converge anche $\sum_n^\infty |c_n w^n|$ converge uniformemente su $|w|\leq\rho<|w_0|$.


Dalla convergenza di $\sum_n c_n w_0^n$ segue che $\lim_n c_nw_0^n=0$ e $|c_nw_0^n|<M$ $\forall n$ quindi \[\sup_{|w|<\rho}|c_n w^n|=\sup_{|w|<\rho} |c_n w_0^n| |\dfrac{w^n}{w_0^n}| \leq M \dfrac{\rho^n}{|w_0^n|}\] con $\rho<|w_0|$ quindi possiamo applicare il criterio di Weierstrass con $a_n=M \dfrac{\rho^n}{|w_0^n|}$ che dà $\sum_n a_n= M \sum_n \dfrac{\rho^n}{|w_0^n|}< \infty$. \qedhere
\end{dimo}

Mettendo insieme queste cose definiamo il raggio di convergenza per una serie di potenze:

\begin{definizione}[Raggio di convergenza]
Per una serie di potenze  $\sum_n c_n (z-a)^n$ chiamiamo $R$ raggio di convergenza se la serie converge uniformemente per $|z-a|\leq R$ e diverge per $|z-a|>R$. 
Sono possibili come valori del raggio di convergenza anche $0$ e $\infty$.
\end{definizione}

Dal teorema di Weierstrass segue che $f(z)=\sum_n^\infty c_n (z-a)^n$ è olomorfa per $|z-a|\leq \rho < R$ dove $R$ è il suo raggio di convergenza.

Ricordiamo brevemente $3$ criteri per ricavare il raggio di convergenza:\begin{itemize}[label=$\cdot$]
\item Criterio del rapporto: Se esiste il limite $\lim_{n}\dfrac{c_n}{c_{n+1}}$ esso coincide col raggio di convergenza
\item Criterio della radice: Se esiste il limite $\lim_{n}\dfrac{1}{\sqrt[n]{|c_n|}}$ esso coincide col raggio di convergenza
\item Criterio di Cauchy-Hadamard: Il raggio di convergenza è $R=\dfrac{1}{\limsup_n \sqrt[n]{|c_n|}}$. 
\end{itemize}

\subsection*{Serie di Taylor}

Sia $f$ olomorfa in $D$ e $\forall z_0\in D$ sia $\eta:=\inf_{w\in \partial D}|w-z_0|$. Allora per ogni $z_0$ la serie di Taylor 
\begin{equation}
\sum_{n=0}^\infty \dfrac{1}{n!} f^{(n)}(z_0)\,(z-z_0)^n
\end{equation}
converge uniformemente a $f(z)$ su ogni disco chiuso centrato in $z_0$ di raggio $\rho$ con $|z-z_0|\leq \rho<\eta$.

\begin{dimo}
Grazie al teorema di Abel ci basterà provare la convergenza puntuale in un qualsiasi $z$ con $|z-z_0|\leq r<\eta$. 

Fissiamo $z$ e scegliamo una circonferenza $\gamma$ centrata in $z_0$ di raggio $|z-z_0|<L<\eta$ qualsiasi, allora dalla formula di Cauchy
\[
f(z)=\dfrac{1}{2\pi i}\oint_\gamma \dfrac{f(w)}{w-z}\, dw
\]
e usando $|w-z_0|=L$ possiamo espandere \[g(w):=\dfrac{f(w)}{w-z}=\dfrac{f(w)}{w-z_0+z_0-z}=\dfrac{f(w)}{w-z_0}\dfrac{1}{1-\dfrac{z-z_0}{w-z_0}}=\dfrac{f(w)}{w-z_0}\sum_{n=0}^\infty \left(\dfrac{z-z_0}{w-z_0}\right)^n\]
dove nell'ultimo passaggio abbiamo riconosciuto una serie geometrica, uniformemente convergente su $\gamma$ per il criterio di Weierstrass con \[a_n=\dfrac{M}{L}\left(\dfrac{|z-z_0|}{L}\right)^n\]
quindi possiamo portarne il segno di serie fuori dall'integrale nella formula di Cauchy ottenendo

\begin{align*}
f(z) &=\dfrac{1}{2\pi i} \oint_\gamma \dfrac{f(w)}{w-z_0} \sum_{n=0}^\infty \left(\dfrac{z-z_0}{w-z_0}\right)^n\, dw =
\dfrac{1}{2\pi i} \sum_{n=0}^\infty \oint_\gamma \dfrac{f(w)}{w-z_0} \left(\dfrac{z-z_0}{w-z_0}\right)^n\, dw =\\
&= \dfrac{1}{2\pi i} \sum_n^\infty (z-z_0)^n \oint_\gamma \dfrac{f(w)}{(w-z_0)^{n+1}} = \quad \text{e dal corollario alla formula di Cauchy} \\ 
&= \sum_n^\infty \dfrac{(z-z_0)^n}{n!}  f^{(n)}(z_0)  \qedhere
\end{align*}

\end{dimo}


\begin{itemize}[label=$\cdot$]
\item Lo sviluppo in serie $f(z)=\sum_n^\infty  c_n (z-a)^n$ è unico.
\item Una funzione che ammette sviluppo di Taylor in un intorno di $z=a$ si dice analitica in $a$.
\item La serie di Taylor di una funzione olomorfa è convergente nel dominio di olomorfia.
\end{itemize}

\begin{esempio}
\begin{itemize}[label=$\cdot$]
\item Lo sviluppo di Mac Laurin di $f(z)=e^z$ è $f(z)=\sum_n^\infty \dfrac{z^n}{n!}$, la funzione è intera e la serie avrà raggio di convergenza $r=\infty$.
\item La funzione $f(z)=\dfrac{1}{1-z}$ è olomorfa in $\mathbb{C}\setminus\{1\}$ quindi possiamo svilupparla in serie di Mac Laurin con raggio di convergenza $R=1$, lo sviluppo è $f(z)=\sum_n^\infty z^n$.
\end{itemize}
\end{esempio}

\subsection*{Serie di Laurent}

Sia $f$ olomorfa in $D$ e in particolare in una corona circolare centrata in $a$ di raggi $0\leq\alpha<\beta\leq\infty$. Allora all'interno della corona	$f(z)$ ammette una rappresentazione in serie
\begin{equation}
f(z)=\sum_{n=-\infty}^{+\infty} c_n (z-a)^n  
\end{equation}
convergente uniformememente in $\alpha'\leq|z-a|\leq\beta'$ con $\alpha'>\alpha$, $\beta'<\beta$.

La successione dei coefficienti è 
\begin{equation}
c_n=\dfrac{1}{2\pi i}\oint_\gamma \dfrac{f(w)}{(w-a)^{n+1}}\,dw
\end{equation}
con $\gamma$ una qualsiasi curva chiusa appartenente alla corona circolare.

\begin{osservazione}
Lo sviluppo di Laurent è esprimibile anche come
\[
f(z)=\sum_{n=0}^{+\infty} c_n (z-a)^n + \sum_{n=1}^{+\infty} \dfrac{d_n}{(z-a)^n} \qquad \begin{array}{c} c_n=\dfrac{1}{2\pi i}\oint_\gamma \dfrac{f(w)}{(w-a)^{n+1}} \\ d_n=\dfrac{1}{2\pi i}\oint_\gamma f(w)(w-a)^{n-1}\, dw\end{array}
\]
dove sono evidenziati il primo termine di Taylor ed il secondo di Laurent.

Chiaramente se la funzione è olomorfa anche nel disco interno alla corona $\gamma$ sarà contraibile in $D$, che sarà anche il dominio di olomorfia di $f(z)(z-a)^{n+1}$ $\forall n\geq 1$, e tutti i coefficienti del termine di Laurent saranno nulli lasciando solo la serie di Taylor.
\end{osservazione}

\begin{dimo}
Possiamo usare la formula di Cauchy con la curva chiusa $\gamma_S=\gamma_\beta+\Gamma-\gamma_\alpha-\Gamma$ completamente contenuta nella corona
\[
f(z)=\dfrac{1}{2\pi i}\int_{\gamma_\beta} \dfrac{f(w)}{w-z}dw - \dfrac{1}{2\pi i} \int_{\gamma_\alpha}\dfrac{f(w)}{w-z}dw
\]

\begin{itemize}
\item Per $w\in\gamma_\beta$ si ha \[\dfrac{f(w)}{w-a-z+a}=\dfrac{f(w)}{w-a}\sum_{n=0}^{+\infty} \left(\dfrac{z-a}{w-a}\right)^n\] con convergenza uniforme in $|z-a|<\beta$
\item Per $w\in\gamma_\alpha$ si ha \[\dfrac{f(w)}{z-a-w+a}=\dfrac{f(w)}{z-a}\sum_{n=0}^{+\infty} \left(\dfrac{w-a}{z-a}\right)^n\] con convergenza uniforme in $|z-a|>\alpha$
\end{itemize}

Quindi per $\alpha<|z-a|<\beta$ entrambe le serie convergono uniformemente e posso portare i segni di sommetoria fuori dagli integrali

\[
f(z)=
\sum_{n=0}^{+\infty}\left[\left(\dfrac{1}{2\pi i}\int_{\gamma_\beta} \dfrac{f(w)}{(w-a)^{n+1}}\, dw\right)\,(z-a)^n\right] - 
\sum_{n=0}^{+\infty}\left[\left(\dfrac{1}{2\pi i}\int_{\gamma_\alpha}f(w) (w-a)^n \,dw\right) \dfrac{1}{(z-a)^{n+1}}\right]
\]

Le due funzioni integrande sono olomorfe e quindi posso deformare a piacere i cammini di integrazione all'interno della corone, scegliendo $\gamma_\beta\sim\gamma_\alpha\sim\gamma$ si trova
\[
f(w)=\sum_{n=0}^{\infty} c_n (z-a)^n + \sum_{n=1}^\infty d_n(z-a)^{n-1} \qquad c_n=\dfrac{1}{2\pi i}\oint-\gamma \dfrac{f(w)}{(w-a)^{n+1}} dw \quad d_n=\dfrac{1}{2\pi i}\oint_\gamma f(w)(w-a)^{n+1} dw
\]
che converge uniformemente in $\alpha'\leq|z-a|\leq\beta'$. \qedhere
\end{dimo}

Lo sviluppo in serie di Laurent è unico.

\begin{esempio}
\begin{itemize}[label=$\cdot$]
\item La funzione $f(z)=e^{\frac{1}{z}}$ è olomorfa in $\mathbb{C}\setminus\{0\}$ e quindi si potrà rappresentare in serie di Laurent attorno all'origine con raggi $\alpha=0$ e $\beta=\infty$. Partendo dallo sviluppo di Taylor di $e^w=\sum_{n=0}^\infty \dfrac{z^w}{n!}$ tramite $w=-z$ troviamo $e^{\frac{1}{z}}\sum_{n=0}^\infty=\dfrac{1}{n!\, z^n}=\sum_{n=-\infty}^{+\infty} \dfrac{z^n}{-n!}$, questa serie converge ad $f$ nella corona e quindi è lo sviluppo in serie di Laurent cercato (per l'unicità).


\item La funzione $f(z)=\dfrac{1}{z-a}$ è olomorfa in $\mathbb{C}\setminus\{a\}$ quindi attorno ad un generico $b\in\mathbb{C}$ avremo due possibili corone in cui svilupparla: $0\leq|z-b|<|a-b|$ e $|z-b|>|a-b|$. Per semplicità sviluppiamo attorno a $b=0$ (corrisponderebbe a ridefinire $w=z-b$ e $a'=a-b$).

La prima corona ha raggio interno $0$ quindi lo sviluppo si riduce alla serie di Taylor $f(z)=\sum_{n=0}^\infty -\dfrac{1}{a} \dfrac{z^n}{a^n}$.

Nella seconda corona manipoliamo la funzione fino a ricondurci ad una serie geometrica $\dfrac{1}{z-a}=\dfrac{1}{\frac{1}{z}-a}=\dfrac{\frac{1}{z}}{1-\frac{a}{z}}=\dfrac{1}{z}\sum_{n=0}^\infty \dfrac{a^n}{z^n}$ e dato che lo sviluppo di Laurent è unico troviamo $f(z)=\sum_{n=1}^\infty \dfrac{a^{n-1}}{z^n}$.
\end{itemize}
\end{esempio}


\subsection{Zeri e poli}

\begin{definizione}[Zero]
Sia $f$ olomorfa in $D$ e $a\in D$. Diciamo che $a$ è uno zero di $f$ di ordine $N$ se $f(a)=0$ e $f(z)=(z-a)^Ng(z)$ con $g$ olomorfa in $D$ e $g(a)\neq0$.\end{definizione}

Chiamiamo semplice uno zero di ordine $N=1$.

\begin{definizione}[Polo]
Sia $f$ olomorfa in $D\setminus\{a\}$ e scegliamo una corona $0<|z-a|<\beta$ con $\beta$ fissato. Allora diciamo che $a$ è un polo di ordine $N$ se nello sviluppo in serie di Laurent di $f$ nella corona ha $c_n=0$ $\forall n<-N$ e $c_{-N}\neq0$. Allora si ha 
\[
f(z)=\dfrac{1}{(z-a)^N}g(z)
\]
con g(z) olomorfa in $0\leq|z-a|<\beta$.
\end{definizione}

Chiaramente se $a$ è un polo di ordine $N$ di $f$ la $(z-a)^N f(z)$ sarà olomorfa in $a$.
Un polo è una singolarità, vedremo meglio più avanti cosa sono le singolarità di una funzione e le proprietà dei poli.

\subsection{Serie di Laurent della reciproca}
f

\section{Singolarità isolate}

\begin{definizione}Un punto $a\in\mathbb{C}$ è una singolarità isolata di una $f$ se esiste un intorno $U\subseteq\mathbb{C}$ di $a$ tale che $f$ sia olomorfa su $U\setminus\{a\}$. Un punto $a\in\mathbb{C}$ è una singolarità non isolata se non ne esiste nessun intorno bucato $U\setminus\{a\}$ in cui $f$ si olomorfa.
\end{definizione}

Una singolarità isolata $z=a$ di una $f(z)$ si dirà:\begin{itemize}[label=$\cdot$]
\item Eliminabile se $|f(z)|<M$ in un intorno $U\setminus\{a\}$ di $a$. 
\item Polare se $\lim_{z\rightarrow a}f(z)=\pm\infty$.
\item Essenziale se non esiste il limite $\lim_{z\rightarrow a}f(z)$.
\end{itemize}

Dallo sviluppo in serie di Laurent attorno ad $a$ di $f(z)=\sum_{n=\infty}^\infty c_n (z-a)^n$ troviamo che la classificazione delle singolarità isolate si può riscrivere come
\begin{itemize}[label=$\cdot$]
\item Eliminabile $\iff$ $c_{-n}=0 \>\> \forall n\in\mathbb{N}\setminus\{0\}$
\item Polare $\iff$ $\exists N\in\mathbb{N}: c_{-n}=0 \>\> \forall n >N$
\item Essenziale $\iff$ la rappresentazione ammette infiniti $c_{-n}$ non nulli
\end{itemize}

\begin{esempio}
\begin{itemize}[label=$\cdot$]
\item La funzione $f(z)=\dfrac{1}{z}$ ha un polo isolato nell'origine.
\item La funzione $f(z)=e^{\frac{1}{z}}$ attorno all'origine ha sviluppo $f(z)=\sum_{n=0}^\infty \dfrac{1}{n!} \dfrac{1}{z^n} = \sum_{n=-\infty}^0 \dfrac{1}{-n!} z^n $
che ha infiniti termini negativi, quindi $z=0$ dovrebbe essere una singolarità essenziale.\par

Vediamo che \[z=|z|e^{i\theta} \implies |f(z)|=\left|e^{\frac{1}{|z|}e^{-i\theta}}\right| = \left|e^{\frac{1}{|z|}(\cos{\theta}-i\sin{\theta})}\right|=e^{\frac{\cos{\theta}}{|z|}}
\]
e quindi il limite dipende da $\theta$ e $\lim_{z\rightarrow o} f(z)$ non esiste.
\end{itemize}
\end{esempio}

\begin{teorema}[Picard]
In ogni intorno bucato di una singolarità essenziale una funzione olomorfa assume tutti i valori in $\mathbb{C}$ tranne al più uno. 
\end{teorema}
\begin{esempio}
Per $f(z)=z$ se consideriamo $f(z_1)=\alpha$ avremo $ \alpha=|\alpha|e^{i\theta}=e^{\frac{1}{z}}$ quindi 
\[
z_{\alpha,k}=\dfrac{1}{\ln{|\alpha|}+i(\theta+2\pi k)}\underset{k\rightarrow \infty}{\longrightarrow}0
\] 
ma non esiste nessuno zero tale che $e^{\frac{1}{z}}=0$.
\end{esempio}

%Parte 2-----------------------------------------------------------------
\cleardoublepage
\addcontentsline{toc}{part}{Spazi di Hilbert e Distribuzioni}
\thispagestyle{plain}
	\begin{center}
		\vspace{.1\textheight}
		{\LARGE\scshape \textsc{Seconda parte}\par}
		\vspace{.1\textheight}
		\textsc{}\par
		\vspace{.05\textheight}
		{\Huge \textsc{Spazi di Hilbert e Distribuzioni}\par}
		\vfill

\end{center}
\fancyfoot[C]{\small Spazi di Hilbert e Distribuzioni}
%------------------------------------------------------------------


\section{Topologie}

Per poter definire i concetti di continuità e convergenza su uno spazio generico è necessaria l'introduzione di una topologia, cioè di una base di intorni.

La gerarchia delle tipologie di spazi che andremo ad introdurre è 
\[
\text{Topologici}\overset{\text{Distanza}}{\longrightarrow}\text{Metrici}\overset{\text{Struttura lineare}}{\longrightarrow}\text{Metrici lineari}\overset{\text{Norma}}{\longrightarrow}\text{Normati}
\overset{\text{Prodotto scalare}}{\longrightarrow}\text{Unitari}
\]


La {\bfseries topologia} $T$ di un insieme $M$ è la collezione dei suoi sottoinsiemi aperti, trovare una topologia di $M$ significa distinguere tra insiemi aperti e chiusi in $M$\footnote{In una definizione rigorosa di topologia si specifica anche la presenza dell'insieme vuoto e la sua chiusura per unioni, anche infinite, e intersezioni finite di aperti.}.

Uno spazio topologico è uno spazio dotato di una topologia.

Consideremo da ora solo spazi topologici.

\begin{definizione}
\begin{enumerate}[i.]
\item Un intorno aperto di $P\in M$ è un insieme aperto $I\subset T$ tale che $P\in I$.
\item Sia $\{P_n\}_n$ una successione im $M$. Diciamo che $\lim_n P_n=P\in M$ se per ogni intorno $I$ di $P$ esiste un $N$ tale che $P_n\in I$ $\forall n>N$.
\item Data una mappa $f:M\rightarrow N$ tra due spazi topologici diciamo che $\lim_{p\rightarrow P_0} f(P)=Q$ se per ogni intorno $I$ di $Q$ esiste un intorno $J$ di $P_0$ tale che $f(J\setminus\{P_0\})\subset I$.
\item Diciamo che una mappa $f:M\rightarrow N$ è continua in $P_0\in M$ se $\lim_{P\rightarrow P_0} f(P)=f(P_0)$
\end{enumerate}
\end{definizione}

\begin{proposizione}
Data una successione $P_n$ in $M$ tale che $\lim_n P_n=P$ ed una mappa $f:M\rightarrow N$ continua in $P$ si ha $\lim_n f(P_n)=f(P)$.
\end{proposizione}

\begin{definizione}[Spazio metrico] Un insieme $A$ si dice spazio metrico se è dotato di una mappa $d:A\times A\rightarrow\mathbb{R}$, chiamata distanza, che $\forall u,v,w\in A$ rispetti\begin{itemize}[label=$\cdot$]
\item $d(u,v)\leq0$ e $d(u,v)=0\iff u=v$
\item $d(u,v)=d(v,u)$
\item $d(u,v)\leq d(u,w)+d(w,v)$
\end{itemize}\end{definizione}

In uno spazio metrico in generale \textbf{non} c'è una struttura lineare, cioè non è uno spazio vettoriale.

\subsection{Basi di intorni}

\begin{definizione}
Dati uno spazio topologico $M$ ed un suo elemento $P\in M$ una famiglia $B(P)$ di aperti di $M$ si chiama una base di intorni di $P$ se:\begin{itemize}[label=$\cdot$]
\item Ogni $I\in B(P)$ è un intorno di $P$
\item Per ogni intorno $J$ di $P$ esiste un $I\in B(P)$ tale che $I\subseteq J$.\end{itemize}
\end{definizione}
	
\begin{esempio}
In $\mathbb{C}^d$ per un generico $z_0$ la famiglia $B(z_0)=\{z: ||z-z_0||<r \text{ al variare di } r\}$ è una base di intorni 
\end{esempio}

Possiamo usare le basi di intorni per definire la continuità di mappe tra spazi topologici
\[f:M\rightarrow N\text{ continua in }x_0\in M \iff \forall\epsilon>0\>\exists\delta>0:|x-x_0|<\delta\implies |f(x)-f(x_0)|<\epsilon \]
La continuità in termini di basi di intorni e quella in termini di intorni sono equivalenti.

Per un punto di uno spazio topologico esiste sempre una base di intorni.

\begin{proposizione}\label{propBaseIntorni}
Sia $M$ uno spazio topologico. Per ogni elemento $P\in M$ sia assegnata una base di intorni $B(P)$. Allora:\begin{itemize}[label=$\cdot$]
\item $B(P)\neq\emptyset$
\item $\forall J\in B(P)$ si ha $P\in J$
\item Dati $J_1,J_2\in B(P)$ $\exists J_3\in B(P)$ tale che $J_3\subset J_1 \cap J_2$
\item Per ogni $Q\in J\in B(P)$ $\exists I\in B(Q)$ tale che $I\subseteq J$
\end{itemize}
\end{proposizione}

\begin{teorema}
Sia $M$ uno spazio topologico e per ogni $P\in M$ sia assegnata una $B(P)$ famiglia di sottoinsiemi di $P\in M$ definita con le proprietà della proposizione \ref{propBaseIntorni}. Allora esiste una e una solo topologia $T$ su $M$ per cui le $B(P)$ sono basi di intorni. Tale topologia è definita da \[A\in T \iff \> \forall P\in A\> \exists J\in B(P) : J\subset A\]
\end{teorema}

\begin{itemize}[label=$\cdot$]
\item La topologia è identificata univocamente tramite la scelta di una base di intorni per ogni punto dello spazio.
\item La continuità può essere definita tramite le basi di intorni.
\item Per spazi vettoriali è sufficiente scegliere una base di intorni per l'origine per identificare la topologia.
\item In uno spazio metrico $M$ con distanza $d$ la famiglia dei $I_\epsilon=\{y\in M : d(x,y)<\epsilon\}\subset B(x)$ è una base di intorni.
\end{itemize}

\section{Spazi di Hilbert}

Ricordiamo alcune definizioni e proprietà note degli spazi vettoriali su campo complesso:
\begin{itemize}[label=$\cdot$]
\item Si dice che uno spazio vettoriale $V$ ha dimensione $n\geq1$ se $V$ contiene $n$ vettori linearmente indipendenti e se ogni insieme di $n+1$ vettori di $V$ è linearmente dipendente.
\item Se uno spazio vettoriale $V$ contiene $n$ vettori linearmente indipendenti per ogni $n\in\mathbb{N}$ si dice che $V$ è $\infty$ dimensionale. 
\end{itemize}

\begin{definizione}[Spazi Unitari o preHilbertiani]
Uno spazio vettoriale $V$ su $\mathbb{C}$ si dice unitario se è dotato di un' operazione $V\times V\rightarrow\mathbb{C}$, detta prodotto scalare, che $\forall u,v,w\in V$, $\forall a\in\mathbb{C}$ soddisfi:\begin{itemize}[label=$\cdot$]
\item $(u,u)\geq 0$ e in particolare $(u,u)=0\iff u=0$
\item $\overline{(u,v)}=(v,u)$
\item $(u,av)=a(u,v)$
\item $(u,v+w)=(u,v)+(u,w)$
\end{itemize}\end{definizione}

\begin{itemize}[label=$\cdot$]
\item Due vettori di uno spazio unitario si dicono ortogonali se hanno prodotto scalare nullo.
\item Se dei vettori $\{v_i\}_{i=1}^n$ di uno spazio unitario sono tutti ortogonali fra loro allora sono un insieme linearmente indipendente.
\item Uno spazio unitario $V$ è anche uno spazio normato con norma indotta dal prodotto scalare $||v||=\sqrt{(v,v)}$ che gode delle usuali proprietà della norma, tra cui $||u+v||\leq ||u||+||v||$
\item In ogni spazio unitario vale la disuguaglianza di Schwarz $|(u,v)|\leq ||u||\, ||v||$ per la norma indotta dal prodotto scalare.
\item La norma indotta rispetta $||u+v||\leq (||u||+||v||)^2$.
\item Uno spazio unitario $V$ è anche uno spazio metrico con metrica $d(u,v)=||u-v||=\sqrt{(u-v,u-v)}$
\item Spazi normati e spazi metrici lineari sono vettoriali quindi sarà sufficiente una base di intorni dell'origine per definire una topologia.
\item In uno spazio normato $V$ la norma $||\cdot||:V\rightarrow\mathbb{C}$ è una mappa continua e in uno spazio unitario $W$ il prodotto scalare $(\cdot,\cdots):W\times W\rightarrow\mathbb{C}$ è una mappa continua.
\end{itemize}

\begin{dimo}[Dimostrazione della disuguaglianza di Schwarz]
\[
0\leq||u-\lambda v||^2=(u-\lambda v,u-\lambda v)=\left| ||u||^2 +|\lambda|^2||v||^2 -\overline{\lambda}(v,u) -\lambda(u,v)  \right|
\] 
e per $\lambda:=\dfrac{(u,v)}{||v||^2}$ si ha 
\[
||u-\lambda v||=||u||^2 + \dfrac{|(u,v)|^2}{||v||^2}-\dfrac{|(u,v)|^2}{||v||^2}-\dfrac{|(u,v)|^2}{||v||^2}\geq 0
\]
quindi $||u||\geq\dfrac{|(u,v)|^2}{||v||^2}$ da cui $||u||^2 \, ||v||^2 \geq |(u,v)|^2$ \qedhere
\end{dimo}


\begin{definizione}[Spazio di Hilbert]
Uno spazio unitario completo per la norma indotta si dice spazio di Hilbert. 

Indicheremo con $\mathcal{H}$ un generico spazio di Hilbert.
\end{definizione}

\begin{itemize}[label=$\cdot$]
\item Uno spazio di Hilbert è anche uno spazio di Banach.
\item Un qualsiasi spazio unitario finito dimensionale è isomorfo a $\mathbb{C}^n$ ed, essendo $\mathbb{C}^n$ completo, è uno spazio di Hilbert.
\end{itemize}

\begin{definizione}
Sia $X$ uno spazio topologico e $S\subset X$. Allora definiamo chiusura di $S$ come 
\[
\overline{S}=\{p\in X \text{ tale che ogni intorno di }p\text{ contenga un punto di }S\}
\]
\begin{itemize}[label=$\cdot$]
\item $S\subseteq X$ si dice chiuso se $S=\overline{S}$ 
\item $S\subseteq X$ si dice denso in $X$ se $\overline{S}=X$
\end{itemize}
\end{definizione}

Se $X$ è uno spazio metrico allora la chiusura è anche $\overline{S}=\{p\in X : \exists\{p_n\}_n\subset S \text{ con } \lim_n p_n=p\}$.

\begin{definizione}\begin{itemize}[label=$\cdot$]
\item Una varietà lineare di uno spazio di Hilbert $\mathcal{H}$ è un suo sottoinsieme $L\subseteq \mathcal{H}$ tale che se $\psi,\chi\in L$ allora per ogni $a,b\in\mathbb{C}$ si ha che $a\psi+b\chi\in L$. In particolare una varietà lineare è uno spazio vettoriale.
\item Una varietà lineare di uno spazio di Hilbert $\mathcal{H}$ si dice un sottospazio di $\mathcal{H}$ se è chiusa.
\end{itemize}
\end{definizione}

\begin{esempio} $C[a,b]$ non è uno spazio di Hilbert perchè non è completo.

 L'insieme $C[a,b]$ è una varietà lineare di $L_2[a,b]$ ma anche se $C[a,b]$ è denso in $L_2[a,b]$ non è chiusa e non è un suo sottospazio.

\end{esempio}

\begin{proposizione}Un sottospazio $S$ di uno spazio di Hilbert è a sua volta uno spazio di Hilbert.
\end{proposizione}

\subsection*{Teorema del completamento}

Per ogni spazio metrico $X$ esiste uno spazio metrico $Y$ per cui: \begin{itemize}[label=$\cdot$] \item $X\subseteq Y$ \item $Y$ è completo \item $X$ è denso in $Y$ \item Per qualsiasi $p,q\in X$ vale $d_Y(p,q)=d_X(p,q)$\end{itemize}

Il completamento è unico a meno di isometrie.

\begin{esempio}
Il completamento di $P[a,b]$ e di $C[a,b]$ è $L_2[a,b]$
\end{esempio}
 
Un qualsiasi spazio vettoriale prehilbertiano $V$ ammetterà dunque un unico completamento $\mathcal{H}$, in cui è denso, che sia uno spazio di Hilbert.



\subsection{Lo spazio delle successioni quadrato sommabili $l_2$\ }

Chiamiamo $l_2$ lo spazio delle successioni $\{u_n\}_{n}$ in $\mathbb{C}$ quadrato sommabili, cioè 

\[l_2=\left\{ u=\{u_n\}_{n=1}^\infty\,, \> u_n\in\mathbb{C} \> \forall n \> : \> \sum_{n=1}^\infty |u_n|^2<\infty\right\}\]

Per costruzione $l_2$ è uno spazio vettoriale, notiamo in particolare che $u,v\in l_2\implies u+v\in l_2$.

In $l_2$ possiamo definire un prodotto scalare con \begin{equation} (z,w)=\sum_{n=1}^\infty \overline{z}_n\,w_n \end{equation} e questo indurrà una norma \begin{equation} ||z||_{l_2}=\sqrt{\sum_{n=1}^\infty |z_n|^2 }\end{equation}

\begin{dimo}[Dimostriamo che $(u,v)<\infty$]
Da $|ab|\leq\dfrac{1}{2}(|a|^2+|b|^2)$ abbiamo $|(u,v)|\leq\sum_{j=1}|u_j|\,|v_j| \leq \dfrac{1}{2}\sum_{j=1}(|u_j|^2+|v_j|^2)=\dfrac{1}{2}\left( ||u||^2+||v||^2\right)<\infty$
\end{dimo}


\begin{teorema}$l_2$ è uno spazio di Hilbert\end{teorema}
\begin{dimo}
Sia $\{u^a\}_{a=1}$ una successione di elementi in $l_2$, cioè ciascun $u^1,\dots,u^a,\dots$ sia una successione quadrato sommabile $\sum_n ||u^a_n||^2<\infty$, e sia $\{u^a\}_a$ una successione di Cauchy.

Allora per ogni $\epsilon>0$ $\exists N$ tale che $a,b>N$ $\implies$ $||u^a-u^b||_{l_2}<\epsilon$ quindi
\[
\sum_{n=1}^\infty \left|u_n^a-u_n^b\right|^2<\epsilon \implies \left|u_n^a-u_n^b\right|<\epsilon \>\forall n \text{ fissato}
\]
Quindi per ogni $n$ fissato la successione $\{u_n^a\}_a$ converge per $a\rightarrow\infty$ (è uniformemente di Cauchy).

Sia $\lim_{a\rightarrow\infty}u_n^a=u_n$, da $\forall\epsilon\exists N:$ per $a,b>N \>\sum_n^\infty |u_n^a-u_n^b|^2<\epsilon$ deduciamo che $\sum_n^M |u_n^a-u_n^b|^2<\epsilon^2$ $\forall M$ e prendendo al limite  $\lim_{b\rightarrow\infty}u_n^b=u_n$ si ha 
\[
\sum_{n=1}^M \left|u_n^a -u_n\right|^2<\epsilon^2
\]
Prendendone il limite per $M\rightarrow\infty$ segue che $||u^a-u||<\epsilon$ e quindi $\lim_{a\rightarrow\infty}u^a=u$.

Dato che $u=\underset{||\,||<\epsilon}{u-u^a}+\underset{u^a\in l_2}{u^a}$ abbiamo che $u\in l_2$.

Quindi ogni successione di Cauchy in $l_2$ converge in $l_2$ quindi $l_2$ è completo ed essendo unitario è uno spazio di Hilbert.\qedhere
\end{dimo}

\subsection{Gli spazi $L_2$\ }

Consideriamo le funzioni $f:[a,b]\rightarrow\mathbb{C}$ e identifichiamo tutte quelle che differiscono solo in sottoinsiemi trascurabili di $[a,b]$ in classi di equivalenza $f_1(x)=f_2(x)\>\text{q.o.}\implies f_1\sim f_2$.

Definiamo $L_2[a,b]$ l'insieme delle classi di equivalenza delle funzioni quadrato integrabili su $[a,b]$
\[
L_2[a,b]=\left\{[f]:\int_a^b |f|^2\,dx<\infty\right\}\qquad\quad [f]=\left\{g:[a,b]\rightarrow\mathbb{C}\text{ tc }g(x)=f(x)\>\text{q.o. in }[a,b]\right\}
\]

$L_2[a,b]$ è uno spazio vettoriale, infatti 
\[
\int_a^b |f+g|^2\,dx=\int_a^b (|f|^2+2|f|\,|g|+|g|^2)\,dx=\int_a^b (|f|^2+|g|^2)\,dx+2(|f|,|g|)<\infty
\] 
quindi $f,g\in L_2[a,b]\implies f+g\in L_2[a,b]$.

Su $L_2[a,b]$ possiamo definire un prodotto scalare 
\begin{equation} 
(f,g)=\int_a^b \overline{f}(x)g(x)\,dx 
\end{equation}
che è ben definito 
\[
|(f,g)|\leq\int_a^b |f|\,|g|\,dx\leq\frac{1}{2}\int_a^b (|f|^2+|g|^2)\,dx<\infty
\]
e da questo definiamo una norma 
\begin{equation}
||f||^2_{L_2}=\int_a^b |f|^2\,dx
\end{equation}

Quindi $L_2[a,b]$ è uno spazio prehilbertiano.

in modo analogo si definiscono $L_2(\mathbb{R})$ e $L_2(A)$ per $A\subseteq\mathbb{R}$ un qualsiasi sottoinsieme misurabile di $\mathbb{R}$, anche di misura infinita.

\[
L_2(A)=\left\{[f]:\int_A |f|^2\,dx<\infty\right\}
\]

\begin{teorema}Gli spazi $L_2(A)$ con $A\subseteq\mathbb{R}^d$ sono completi e quindi di Hilbert\end{teorema}

\begin{osservazione}
Tutti gli spazi $L_p(A)=\{f:A\rightarrow\mathbb{C} : \int_A |f|^p\,dx<\infty\}$ con $1\leq p<\infty$ sono spazi di Banach con norma $||f||_p=\left(\int_A |f|^p \, dx\right)^{\frac{1}{p}}$.

L'unico $L_p$ che è anche uno spazio di Hilbert però è $L_2$ perchè per gli altri non è ben definito un prodotto scalare.


Definiamo inoltre \[L^\infty(A)=\{f:A\rightarrow\mathbb{C} : ||f||_\infty<\infty\}\] dove $||f||_\infty=\sup_{x\in A}|f(x)|$ e notiamo che se $f\in L^p\cap L^\infty$ $\forall p\geq 1$ allora \\$||f||_\infty = \lim_{p\rightarrow\infty} ||f||_p$
\end{osservazione}

\subsection{Decomposizione in spazi ortogonali}

\begin{definizione}
Dato un sottoinsieme $S\subset\mathcal{H}$ definiamo il complemento ortogonale 
\[
S^\perp=\left\{ \psi\in\mathcal{H} : (\psi,\chi)=0 \> \forall\chi\in S\right\}
\]
\end{definizione}

Dalla definizione si ha che $S^\perp$ è una varietà lineare e che è chiuso per la continuità del prodotto scalare $(\cdot,\chi)$. Quindi $S^\perp$ è un sottospazio di $\mathcal{H}$ ed è uno spazio di Hilbert.

\begin{teorema}
Sia $S$ un sottospazio di $\mathcal{H}$. Allora $\psi_\perp\in S^\perp$ e $\mathcal{H}=S\oplus S^\perp$ e ogni $\psi\in\mathcal{H}$ ammette un'unica decomposizione $\psi=\psi_\parallel+\psi_\perp$ con $\psi_\parallel\in S$. 
\end{teorema}

\begin{dimo}[Dimostrazione dell'unicità]
Sia $\psi=\psi_\parallel^1+\psi_\perp^1=\psi_\parallel^2+\psi_\perp^2$ 

Allora $\psi_\parallel^1-\psi_\parallel^2=\psi_\perp^2-\psi_\perp^1$ ma quindi $\left(\psi_\parallel^1-\psi_\parallel^2,\psi_\perp^2-\psi_\perp^1\right)=0=\left(\psi_\perp^2-\psi_\perp^1,\psi_\perp^2-\psi_\perp^1\right)=||\psi_\perp^2-\psi_\perp^1||^2=||\psi_\parallel^1-\psi_\parallel^2||^2=0$ quindi $\psi_\perp^1=\psi_\perp^2$ e $\psi_\parallel^1=\psi_\parallel^2$.\qedhere 

\end{dimo}

Come conseguenza di questo teorema se $S$ è una varietà lineare in $\mathcal{H}$ si ha che\begin{itemize}[label=$\cdot$]\item $\overline{S}^\perp=S^\perp$ \item $\mathcal{H}=\overline{S}\oplus S^\perp$ \item $\left(S^\perp\right)^\perp=\overline{S}$ \item $\overline{S}=\mathcal{H} \iff S^\perp=\emptyset $\end{itemize}

\subsection{Basi ortonormali}

In uno spazio vettoriale finito dimensionale $V$ con $\dim V=N$ una base ortonormale è un insieme $\{e_i\}_{i=1}^N$ con $(e_i,e_j)=\delta_{ij}$ tale che per ogni $v,w\in V$:\begin{itemize}[label=$\cdot$]
\item $v=\sum_j v_j e_j$ con $v_j=(e_j,v)$
\item $||v||=\sqrt{\sum_j |v_j|^2}$
\item $(v,w)=\sum_j \overline{v}_j w_j$
\end{itemize}

\begin{definizione}Un sistema ortonormale di vettori di uno spazio prehilbertiano $V$ è una collezione $\{x_i\}_i$ di vettori $x_i\in V$ tali che $(x_i,x_j)=\delta_{ij}$.
\end{definizione}


\begin{teorema}[Disuguaglianza di Bessel] Sia $\{x_i\}_i^N$ un sistema ortonormale di uno spazio prehilbertiano $V$. Allora per un $y\in V$ \[\left|\left| y\right|\right|^2\geq\sum_{n=1}^N \left|(y,x_n)\right|^2\]
\end{teorema}

\begin{dimo} 
$y=\psi+\phi$ con $\psi=\sum_i^n (x_i,y)x_i$ e $\phi=y-\sum_i^n (x_i,y)x_i$. Allora

\begin{align*} 
(\psi,\phi)&=\sum_{n=1}^N \overline{(x_n,y)}(x_n,y)-\sum_{n,k=1}^N \overline{(x_n,y)} \left(x_n,(x_k,y)x_k\right)=\\
&==\sum_{n=1}^N \overline{(x_n,y)}(x_n,y)-\sum_{n,k=1}^N \overline{(x_n,y)} (x_k,y)   \,
\underset{=\delta_{nk}}{\left(x_n,x_k\right)}
=0
\end{align*}
Quindi
\begin{align*} 
||y||^2&= ||\psi||^2+||\phi||^2+2(\psi,\phi)\geq||\psi||^2=\sum_{n=1}^N \left|(y,x_n)\right|^2 \qedhere
\end{align*}
\end{dimo}

\begin{definizione}Una base ortonormale di uno spazio di Hilbert $\mathcal{H}$ è un sistema ortonormale $S$ di $\mathcal{H}$ tale che nessun altro sistema ortonormale di $\mathcal{H}$ lo contenga propriamente.\end{definizione}

\begin{teorema}Ogni spazio di Hilbert ammette una base ortonormale\end{teorema}

Omettiamo la dimostrazione di questo teorema, accenniamo soltanto che usa il Lemma di Zorn e segue la stessa idea di quella dell'assioma della scelta.

\begin{lemma}\label{lemmaSistemaBaseDenso}
Un sistema ortonormale $S=\{e_\alpha\}_{\alpha\in A}$ di $\mathcal{H}$ è una base ortonormale di $\mathcal{H}$ se e solo se l'insieme $M$ di tutte le combinazioni lineari finite di $S$ è denso in $\mathcal{H}$.
\end{lemma}

\begin{dimo}
$\overline{M}=\mathcal{H}\implies M^\perp=\emptyset$  quindi non esiste nessun vettore ortogonale a $S$ quindi $S$ è una base.

\end{dimo}

\begin{definizione}Uno spazio di Hilbert $\mathcal{H}$ si dice separabile se ne esiste un sottoinsieme $W$ numerabile e denso in $\mathcal{H}$\end{definizione}

\begin{teorema}Uno spazio di Hilbert è separabile se e solo se ammette una base ortonormale numerabile\end{teorema}

\begin{dimo}~

\begin{itemize} 
\item Sia $\mathcal{H}$ separabile con $W$ un sottoinsieme numerabile denso $\overline{W}=\mathcal{H}$

Eliminiamo da $W$ i vettori linearmente dipendenti da altri in modo da ottenere un insieme $Z\subseteq W$ numerabile e linearmente indipendente il cui insieme delle combinazioni lineari (appunto $W$) sia denso in $\mathcal{H}$. Allora $Z=\{e_n\}_n^\infty$. Col procedimento di Gram-Schmidt otteniamo poi un altro insieme $S=\{y_n\}_n^\infty$ con $(y_n,y_m)=\delta_{nm}$ il cui insieme delle combinazioni lineari è ancora denso in $\mathcal{H}$. Quindi per il lemma \ref{lemmaSistemaBaseDenso} $S$ è una base.
\item Se $\mathcal{H}$ ammette una base numerabile allora ne esiste un sottoinsieme numerabile denso, infatti per $y\in\mathcal{H}$ è sufficiente considerare ``l'approssimante'' $y^*=\sum_{j=1}^N c_j^* e_j$ con $c_j$ i numeri razionali che approssimano $(e_j,y)$.\qedhere 
\end{itemize}
\end{dimo}

\begin{teorema}
Sia $\mathcal{H}$ uno spazio di Hilbert separabile e sia $S=\{e_n\}_n^\infty$ una base ortonormale numerabile.
Allora valgono:\begin{enumerate}[i.]
\item Ogni elemento di $\mathcal{H}$ ammette l'espansione in serie di Fourier $y=\sum_n^\infty (e_n,y)e_n$ 
\item \textbf{ L'indentità di Parseval} \[||y||^2=\sum_n^\infty \left|(e_n,y)\right|^2 \qquad \quad(x,y)=\sum_n^\infty \overline{(e_n,x)}(e_n,y)\]
\item Il \textbf{teorema di Riesz-Fischer} 
\[ \sum_n^\infty|c_n|^2<\infty \implies \sum_n^\infty c_n e_n\in\mathcal{H}\]
\end{enumerate}
\end{teorema}

\begin{dimo}


\end{dimo}

Questo ci dice che tutti gli spazi di Hilbert separabili sono isomorfi a $l_2$.

\begin{teorema}
Tutti gli spazi $L_2(A)$ con $A\subseteq\mathbb{R}^d$ sono separabili.
\end{teorema}

\begin{esempio}[BasI di $L_2(A)$]
\begin{itemize}[label=$\cdot$]
\item Base di Fourier $L_2[0,L]$

\item Base di Legendre $L_2[-1,1]$

\item Base di Hermite $L_2(\mathbb{R})$

\end{itemize}
\end{esempio}

\section{Funzionali lineari}

Un funzionale è una mappa $V\rightarrow\mathbb{C}$ dove $V$ può essere uno spazio arbitrario qualsiasi.

\begin{esempio}
Scegliendo $V=C[a,b]$ la mappa $f\mapsto\int_a^b |f(x)\,dx$ è un funzionale.
\end{esempio}

\begin{definizione}
Un funzionale lineare su uno spazio vettoriale $U$ è una mappa $F:U\rightarrow\mathbb{C}$ lineare
\[
F(au+bv)=aF(u)+bF(v)\qquad\quad\forall u,v\in U\> \forall a,b\in\mathbb{C}
\]
\end{definizione}

\begin{esempio}~

Su $\mathbb{C}$ la mappa $F(u)=ku$ con $k\in\mathbb{C}$ fissato è un funzionale lineare.

Su $\mathbb{C}^n$ la mappa $F(u)=(w,u)$ con $w\in\mathbb{c}^n$ fissato è un funzionale lineare.
\end{esempio}

\begin{definizione}
Una mappa lineare $F:U\rightarrow V$ con $U,V$ spazi normati si dice limitata se esiste un $M>0$ tale che 
\[
\left|\left| F(u)\right|\right|_V\leq M\left|\left|u\right|\right|_U
\]
per ogni $u\in U$.

Il più piccolo $M$ per cui vale tale proprietà è chiamato norma di $F$ ed indicato con $||F||$. 
\end{definizione}
Chiaramaente la norma di $F$ così definita rispetta
\begin{equation}
\left|\left|F\right|\right|=\sup_{u\in U} \dfrac{\left|\left|F(u)\right|\right|_V}{||u||_U}=\sup_{||u||_U=1} \left|\left|F(u)\right|\right|_V
\end{equation}


\begin{teorema}
Una mappa lineare $F$ tra spazi normati $U,V$ è limitata se e solo se è continua.
\end{teorema}
\begin{dimo}~
\begin{enumerate}[i.]
\item Sia $F$ limitata. Per mostrare la continuità occorre dimostrare che $\lim_{h\rightarrow0}F(u+h)=F(u)$ con $u,h\in U$, ma $\left|\left|F(u+h)-F(u)\right|\right|=\left|\left|F(h)\right|\right|\leq M ||h||$ e quindi a limite si ha la continuità.
\item Sia $F$ continua, allora esiste $\Delta>0$ tale che 
\[
||h||\leq\Delta \implies \left|\left| F(h) \right|\right| = \left|\left| F(h)-F(0) \right|\right|\leq 1 
\]
per $u\in U$ quindi
\[
\left|\left| F(u)\right|\right| =\left|\left|F\left(\dfrac{u}{||u||}\Delta\right)\right|\right|\, \dfrac{||u||}{\Delta}
\]
ma $\left|\left|\dfrac{u}{||u||}\Delta\right|\right|=||\Delta||$ quindi $\left|\left|F(u)\right|\right|\leq \dfrac{||u||}{\Delta}$ e quindi $F$ è limitato.\qedhere
\end{enumerate}
\end{dimo}

 Lo spazio delle mappe lineari e continue tra due $U,V$ normati è a sua volta uno spazio normato con norma $||F||=\sup_{||u||_U=1}||F(u)||_V$, inoltre se $U,V$ sono spazi di Banach anche lo spazio delle mappe lineari continue è uno spazio di Banach.

In particolare se $U$ è uno spazio di Banach e $V=\mathbb{C}$ l'insieme dei funzionali lineari e continui su $U$ è uno spazio di Banach chiamato \textbf{\boldmath spazio duale di $U$} ed indicato con $U^*$.
La norma su $U^*$ è \begin{equation} ||F||_{U^*}=\sup_{u\in U}\dfrac{|F(u)|}{||u||_U}\end{equation} dove chiaramente $||F(u)||_{\mathbb{C}}=|F(u)|$.

Ora vogliamo individuare lo spazio duale ad un generico spazio di Hilbert $\mathcal{H}$. Per un $w\in\mathcal{H}$ il funzionale lineare $F:\mathcal{H}\rightarrow\mathbb{C}$ definita da $x\mapsto(w,x)$ per un $w\in\mathcal{H}$ fissato è continua e ha norma $||F||_{\mathcal{H}^*}=||w||_{\mathcal{H}}$.

Infatti
\[
\left|F(x)\right|=|(w,x)|\leq ||w|| \,||x||=M||x|| \qquad M=||w||
\]

\begin{teorema}[Teorema di Riesz]
Per ogni $F\in\mathcal{H}^*$ esiste un unico $w\in\mathcal{H}$ tale che $F(x)=(w,x)$ per ogni $x\in\mathcal{H}$, ovvero $\mathcal{H}^*=\mathcal{H}$.
\end{teorema}

\begin{dimo}

\end{dimo}


\begin{esempio}
Il funzionale lineare su $L_2[a,b]$ definito da 
\[
F(\varphi)=\int_a^b \varphi(x)\,dx \qquad\quad \varphi\in L_2[a,b]
\]
è limitato, infatti per $w=1\in L_2[a,b]$ si ha che $F(\varphi)=(w,\varphi)_{L_2[a,b]}$ e quindi \\$||F||_{L_2[a,b]}=||w||_{L_2[a,b]}=\sqrt{b-a}$.
\end{esempio}

\section{Distribuzioni}

Vogliamo allargare il concetto di ``funzione'' introducendo le distribuzioni:\begin{itemize}
\item Le distribuzioni ammettono sempre derivata
\item Le funzioni sono distribuzioni
\item Le distribuzioni ammettono sempre trasformata di Fourier
\item Le distribuzioni permettono la risouluzioni delle equazioni differenziali alle derivate parziali
\end{itemize}



\subsection{Spazio delle funzioni di test}

Lo spazio delle funzioni di Test di Schwartz è definito da
\begin{equation}
S(\mathbb{R})=\left\{\varphi\in C^\infty(\mathbb{R}) : \sup_{x\in\mathbb{R}}\left|x^i D^j \varphi\right|=:||\varphi||_{ij}<\infty \>\forall i,j\in\mathbb{N}\right\}
\end{equation}

le funzioni di test sono funzioni di classe $C^\infty$ le cui derivate di ogni ordine decrescono ad infinito più velocemente dell'inverso di qualsiasi potenza.

\begin{esempio}~
\begin{itemize}[label=$\cdot$]
\item $\dfrac{1}{x^2+1}\not\in S(\mathbb{R})$ perchè $\dfrac{x^3}{x^2+1}\underset{x\rightarrow\infty}{\longrightarrow}\infty$
\item $e^{-|x|}\not\in S(\mathbb{R})$ perchè non è $C^\infty$
\item $x^me^{-x^2}\in S(\mathbb{R})$: è di classe $C^\infty$ e a limite $x^{m+n}e^{-x^2}\underset{x\rightarrow\infty}{\longrightarrow}0$ $\forall n$ e per ogni ordine di derivata.
\end{itemize}
\end{esempio}

Similmente possiamo definire $S(\mathbb{R}^2)$ con
\[
S(\mathbb{R}^2)=\left\{ \varphi:\mathbb{R}^2\rightarrow\mathbb{C} \> \text{ di classe } C^\infty : \sup_{(x,y)\in\mathbb{R}^2}\left|x^my^n\partial_x^i\partial_y^j \varphi(x,y)\right|=:||\varphi||_{\underset{nj}{mi}}<\infty\>\forall m,n,i,j\in\mathbb{N}\right\}
\]

e allo stesso modo ogni altro $S(\mathbb{R}^d)$.

Lo spazio delle funzioni di test è uno spazio vettoriale infinito dimensionale, in particolare $S(\mathbb{R}^d)\subset L_2(\mathbb{R}^d)$. 


$S(\mathbb{R}^d)$ è denso in $L_2(\mathbb{R}^d)$ per la norma di $L_2$ cioè
\[
\left.\overline{S\left(\mathbb{R}^d\right)}\right._{||\cdot||_2}=L_2\left(\mathbb{R}^d\right)
\]
infatti $L_2(\mathbb{R}^d)$ ammette la base $\{e_n(x)\}_n$ di Hermite con $e_n(x)\in S(\mathbb{R})$.

\begin{definizione}[Topologia di $S$]
Definiamo la base di intorni $B(0)$ di $\varphi(x)=0$, origine di $S(\mathbb{R})$, come le intersezioni finite 
degli insiemi 
\[
U_{ij,\epsilon}=\left\{\phi\in S(\mathbb{R}) : ||\phi||_{ij}<\epsilon\right\}
\]
al variare di $i,j$ ed $\epsilon$.
\end{definizione}

Da questo segue che $S-\lim_{n\rightarrow\infty}\varphi_n=\varphi$ se e solo se $\lim_n ||\varphi_n-\varphi||_{ij}=0$ $\forall i,j$.

\begin{osservazione}[Nota]
\[
\lim_n ||\varphi_n-\varphi||_{ij}=0 \iff \forall\epsilon>0\> \exists N : n>N \implies(\varphi_n-\varphi)\in U_{ij,\epsilon} \text{ cioè }||\varphi_n-\varphi||_{ij}<\epsilon 
\]


In realtà $S(\mathbb{R})$ è metrizzabile con una metrica che induce questa stessa topologia, però la metrica è abbastanza complicata.
\end{osservazione}

La topologia di $S$ è costruita in modo che $S-\lim_n \varphi_n = \varphi$ implichi $L_2-\lim_n \varphi_n=\varphi$ ma non vale il viceversa.

\begin{teorema}$S(\mathbb{R})$ è uno spazio completo.\end{teorema}

\subsubsection{Mappe continue in \boldmath$S(\mathbb{R})$}
\begin{itemize}
\item La mappa $D:S(\mathbb{R})\rightarrow S(\mathbb{R})$ definita da $\varphi\mapsto\varphi'$ è una mappa continua. Fissato un intorno $\tilde{U}$ di $\varphi'=0$ esiste un intorno di $\varphi=0$ tale che $D\phi=\phi'\in\tilde{U}$ se $\phi\in U$. Infatti per come abbiamo definito le seminorme segue che 
\[
\left|\left| D\phi\right|\right|_{ij}=||\phi||_{i,j+1}<\epsilon
\] 

\item Per funzioni $f$ sufficientemente ``piccole'' la mappa data dal prodotto $f\varphi$ è una mappa continua in $S(\mathbb{R})$
\begin{definizione}Definiamo 
$O_M=\left\{f\in C^\infty(\mathbb{R}) : \left|D^i f(x)\right|\leq C_i(1+x^{2n_i})\>\forall i\in\mathbb{N}\right\}$
\end{definizione}
e allora la mappa $G:S(\mathbb{R})\rightarrow S(\mathbb{R})$ data da $\varphi\mapsto f\varphi$ con $f\in O_M$ è continua, infatti senza andare nei dettagli dei conti
\[
\left|\left|f\varphi\right|\right|_{ij}=\sup_{x\in\mathbb{R}} \left|x^iD^j (f\varphi)\right| \leq  \sum_{i,j}^N C_{ij} ||\varphi||_{ij}
\]
\item La mappa $\varphi_0 * :S(\mathbb{R})\rightarrow S(\mathbb{R})$, data dalla convoluzione con una $\varphi_0\in S(\mathbb{R})$ fissata, è una funzione continua in $S(\mathbb{R})$.
\[
\phi\mapsto \varphi_0 * \phi = \int \varphi_0(x-y) \phi(y)\, dy
\]


\end{itemize}

\subsection{Spazio delle distribuzioni temperate}


\begin{definizione}Lo spazio duale a $S(\mathbb{R})$, indicato con $S^*(\mathbb{R})$ o $S'(\mathbb{R})$, si chiama spazio delle distribuzioni temperate o spazio delle distribuzioni di Schwartz.
\[
S^*(\mathbb{R})=\left\{F :S(\mathbb{R})\rightarrow\mathbb{C}, \> \varphi\mapsto F(\varphi)\in\mathbb{C}  \> \text{ lineare, continuo e limitato}\right\}
\]
\end{definizione}

$S'$ è lo spazio vettoriale (per definizione di spazio duale) dei funzionali lineari su $S$ dunque per ogni intorno $\epsilon$ dell'origine in $\mathbb{C}$ esiste un intorno base $U$ di $\varphi=0$ in $S$ tale che se $\varphi\in U$ allora $|F(\varphi)|<\epsilon$.

\begin{teorema}
Un funzionale lineare $F:S(\mathbb{R})\rightarrow\mathbb{C}$ appartiene a $S'(\mathbb{R})$ se e solo se esiste una costante $C>0$ e un insieme finito di seminorme $||\cdot||_{ij}$ con $i,j\in I$ tali che 
\[
\left|F(\varphi)\right|\leq C\sum_{ij\in I} ||\varphi||_{ij} \qquad \forall \varphi\in S(\mathbb{R})
\]
\end{teorema}

\begin{dimo}[Dimostriamo solo il $\Leftarrow$]
Fissato $\epsilon>0$ sia $N$ il numero delle seminorme che esiste per l'ipotesi, scelgo come intorno base $||\varphi||_{ij}<\dfrac{\epsilon}{NC}$ $\forall i,j\in I$. Allora 
\[
|F(\varphi)|\leq C\sum_{ij\in I} \dfrac{\epsilon}{NC}=\epsilon \qedhere
\]
\end{dimo}

Alcuni elementi di $S'$ sono rappresentati da ``vere'' funzione $f:\mathbb{R}^D\rightarrow\mathbb{C}$ e corrispondono ai funzionali definiti da
\[
F(\varphi)=\int_{\mathbb{R}^D} f(x)\varphi(x) \, d^Dx \qquad \varphi\in S
\]

le distribuzioni temperate di questo tipo sono chiamate \textbf{distribuzioni regolari}.

In particolare si ha:\begin{itemize}
\item $S\subset S'$
Sia $\varphi_0\in S$ e poniamo $F(\varphi)=\int_{\mathbb{R}^D} \varphi_0(x)\varphi(x)\,d^Dx$ allora 
\[
|F(\varphi)|\leq\int_{\mathbb{R}^D} |\varphi_0(x)|\,|\varphi(x)|\,d^Dx \leq ||\varphi||_{00} \, \int_{\mathbb{R}^D} |\varphi_0|\,d^Dx= ||\varphi_0||_{L_1}\,||\varphi||_00
\]
\item $L^p\subset S'$
Sia $f\in L^p(\mathbb{R}^D)$ e $\begin{cases}A=\{x\in\mathbb{R}^D : |f|\leq1\} \\ B=\{x\in\mathbb{R}^D : |f|>1\}\end{cases}$, poniamo $F(\varphi)=\int_{\mathbb{R}^D} f(x)\varphi(x)\,d^Dx$ allora
\begin{align*}
|F(\varphi)| & \leq \int_A |f(x)| \,|\varphi(x)|\,d^Dx + \int_B |f(x)| \,|\varphi(x)|\,d^Dx \leq 
\int_A |\varphi(x)|\,d^Dx + \int_B |f(x)|^p \,|\varphi(x)|\,d^Dx \\
& \leq \int_{\mathbb{R}^D} \dfrac{1+|x|^{2D}}{1+|x|^2}\,|\varphi(x)|\,d^D x + \sup_{x\in\mathbb{R}^D}|\varphi(x)|\,\int_{\mathbb{R}^D} |f(x)| \,d^Dx = \int_{\mathbb{R}^D} \dfrac{1+|x|^{2D}}{1+|x|^2}\,|\varphi(x)|\,d^D x + ||f||_{L^p}\,||\varphi||_{00} \\
& \leq \left(\int_{\mathbb{R}^D} \dfrac{d^Dx}{1+|x|^{2D}} \right)\sup_{x\in\mathbb{R}^D}\left|(1+x^{2D})\varphi\right| + a ||\varphi||_{00} =\\
&= b\left(||\varphi||_{00} +\sum_{\{j\}} ||\varphi||_{j_1,...,j_{2d}} \right) + a ||\varphi||_{00}
\end{align*}

\item $O_M=O_M=\left\{f\in C^\infty(\mathbb{R}) : \left|D^i f(x)\right|\leq C_i(1+x^{2n_i})\>\forall i\in\mathbb{N}\right\}\subset S'$

\item Tutte le funzioni localmente integrabili, $\int_{|x|<R}|f(x)dx<\infty$, e polinomialmente limitate all'infinito, $|f(x)|<C(x^2)^N$ per $|x|>R$, sono in $S'$
\begin{align*} 
|F(\varphi)|&\leq \int_{|x|<R} |f|\,|\varphi|\,d^Dx + \int_{|x|>R}|f|\,|\varphi|\,d^Dx \\ 
&\leq ||f||_{\underset{|x|<R}{L^1}}\,||\varphi||_{00} + C\int_{|x|>R}(x^2)^N |\varphi(x)| \, d^Dx \\
& = a ||\varphi||_{00} + C\int_{|x|>R} \dfrac{(x^2)^N}{(x^2)^{N+D}}(x^2)^{N+D}|\varphi|\,d^Dx \\ 
& \leq a||\varphi||_{00} +C\int{|x|>R} \dfrac{d^Dx}{|x|^{2D}} \sum_{\{j\}} c_{\{j\}} ||\varphi||_{j_1,...,j_{2N+2D}}
\end{align*}
\end{itemize}

Per esempio:\begin{itemize}[label=$\cdot$]
\item $x^\alpha$ per $\alpha>-1$ è in $S'$
\item $\left(\ln{|x|}\right)^\beta$ per $\beta>0$ è in $S'$
\item La funzione di Heaviside è una distribuzione temperata
\item $e^{-|x|}\in S'$
\end{itemize}
mentre alcune funzioni ``usuali'' non sono distribuzioni: $e^x\not\in S'$ e $\frac{1}{x}\not\in S'(\mathbb{R})$ per esempio.


\subsection{Distribuzioni non regolari}

In generale le distribuzioni \textbf{non} sono rappresentabili con funzioni. La notazione ``formale'' $F(\varphi)=\int F(x)\varphi(x)\,dx$ si usa spesso nonostante non esista una funzione $F(x)$ per cui sia vera l'uguaglianza:\begin{itemize}[label=$\cdot$]
\item $\delta$ di Dirac piccata in $x=a$ è definita con
\begin{equation}
\delta_a (\varphi):= \varphi(a) \quad \text{ formalmente } \quad \int_{\mathbb{R}^D} \delta^D(x-a) \varphi(x)\,d^Dx=\varphi(a)
\end{equation}
\item La parte principale di $\frac{1}{x}$ è una distribuzione: sappiamo già che $\frac{1}{x}\not\in S'(\mathbb{R})$ perchè non è localmente integrabile. Poniamo
\[
\left(\mathcal{P}\left(\sfrac{1}{x}\right)\right)(\varphi):=\lim_{\epsilon\rightarrow0}\int_{|x|>\epsilon}\dfrac{1}{x}\varphi(x)\,dx = \lim_{\epsilon\rightarrow0} \int_\epsilon^\infty \underset{\text{continua in }x=0}{\dfrac{\varphi(x)-\varphi(-x)}{x}}\,dx = \int_0^\infty \dfrac{\varphi(x)-\varphi(-x)}{x}\,dx
\]
e quindi definiamo il funzionale come
\begin{equation}
\left(\mathcal{P}(\sfrac{1}{x})\right)(\varphi)=\int_0^\infty \dfrac{\varphi(x)-\varphi(-x)}{x}\,dx
\end{equation}
Chiaramente $\mathcal{P}(\sfrac{1}{x})$ è lineare, inoltre
\begin{align*}
\left|\dfrac{\varphi(x)-\varphi(-x)}{x}\right|&=\left|\frac{1}{x}\int_{-x}^x \varphi'(y)\,dy \right| \leq \frac{1}{|x|}\int_{-x}^x \, \underset{\text{costante}}{||\varphi||_{01}}\,dx \leq 2 ||\varphi||_{01}\\
(\mathcal{P}(\sfrac{1}{x}))(x)&\leq \int_0^1 \left|\dfrac{\varphi(x)-\varphi(-x)}{x}\right|\,dx + \int_1^\infty \left|\dfrac{\varphi(x)-\varphi(-x)}{x}\right|\,dx\leq \\ &\leq 2||\varphi||_{01} + \int_1^\infty \left(x|\varphi(x)|+x|\varphi(-x)|\right)\,\dfrac{dx}{x^2}\leq 2||\varphi||_{01}+2||\varphi||_{10}
\end{align*}
quindi $\mathcal{P}(\sfrac{1}{x})$ è una distribuzione.
\end{itemize}

In generale quindi una distribuzione non ``assume valori'' e non si ``annulla'' in un punto $x\in\mathbb{R}^D$, possiamo però specificare cosa intendiamo per annullamente in un aperto $A\subset\mathbb{R}^d$.
\begin{definizione}
Sia $A$ un aperto di $\mathbb{R}^D$ e $F\in S'(\mathbb{R}^D)$. Diremo che $F$ è nulla in $A$ se $F(\varphi)\equiv 0_{S(\mathbb{R}^D)}$ per ogni $\varphi$ il cui supporto sia contenuto in $A$.

Definiamo supporto di $F$, $\text{supp}\,F$, il complemento del più grande $A$ sul quale $F$ è nulla.
\end{definizione}
Pertanto $\text{supp}\,F$ è chiuso.

Per esempio:
\begin{itemize}[label=$\cdot$]
\item $\text{supp}\,\delta_a=\{a\}$
\item $\text{supp}\,H=\mathbb{R}^+$
\end{itemize} 

\subsection{Topologia debole di $S'$}

Dato che $S'$ è uno spazio vettoriale ci basta edfinirne una base di intorni di $F=0$. 
\begin{definizione}
La base di intorni in $S'$ di $F=0$ è data dalle \textbf{intersezioni finite} degli insiemi
\[
U_{\varphi,\epsilon}:=\left\{F\in S' : |F(\varphi)<\epsilon\right\} \qquad \text{ al variare di }\varphi,\epsilon
\]
\end{definizione}

Da questo segue
\begin{equation}
S'\lim_{n\rightarrow\infty}F_n=F\iff\lim_{n\rightarrow\infty}F_n(\varphi)=F(\varphi)\in\mathbb{C} \qquad \forall\varphi\in S
\end{equation}

\begin{itemize}[label=$\cdot$]
\item La topologia indotta dalla base di intorni di $S'$ che abbiamo definito si chiama \textbf{topologia debole} di $S'$.
\item Dato che $L_2\subset S'$ possiamo usare la topologia debole anche in $L^2$, questa sarà però più debole della topologia in norma che abbiamo definito precedentemente per $L^2$, cioè ci saranno successioni convergenti in $L^2$ per la topologia debole ma non per la topologia in norma.
\end{itemize}

Abbiamo inoltre in criterio sufficiente alla continuità in $S'$ di una mappa $A:S'\rightarrow S'$ 
\[
|AF(\varphi)|\leq|F(\varphi_0)|\implies A \text{ continua perchè } \epsilon>k|F(\varphi_0)|\geq|AF(\varphi)|\rightarrow \delta=\dfrac{\epsilon}{k}
\]

\begin{teorema}
S è completo nella topologia debole e $\overline{S}_{debole}=S'$.
\end{teorema}
Questo teorema ci dice che ogni distribuzione temperata è approssimabile nella topologia debole tramite funzioni di test.


Non dimostriamo questo teorema ed il lemma successivo.
\begin{lemma}[Convergenza dominata]
Sia $f_n\in L^1(\mathbb{R}^D)$ ed esista $\lim_{n\rightarrow\infty}f_n(x)=f(x)$ quasi ovunque. Esista inoltre $g(x)\in L^1(\mathbb{R}^D)$ tale che $|f(x)\leq g(x)$ q.o., allora
\[
\lim_{n\rightarrow\infty}\int f_n(x)d^Dx=\int \lim_{n\rightarrow\infty}f_n(x)d^Dx=\int f(x) d^Dx
\]
\end{lemma}

\subsubsection{Serie \boldmath$\delta$}

Le successioni $\delta$ sono successioni di funzioni che convergono, per la topologia debole, alla delta di Dirac.

Queste successioni si possono definire in modo equivalente come $F_n$ convergenti alla $\delta$ per $n\rightarrow\infty$ o per $\epsilon=\frac{1}{n}$ che convergeranno alla $\delta$ per $\epsilon\rightarrow0$.

\begin{enumerate}[i.]
\item $F_n(x)=\dfrac{n}{\sqrt{\pi}}e^{-n^2x^2}\in S(\mathbb{R})$
\item $F_n(x)=\dfrac{\sin{nx}}{\pi x}\in L^p(\mathbb{R})$ con $p\geq2$
\item $F_n(x)=\dfrac{n}{\pi(n^2x^2+1)}=\dfrac{\epsilon}{\pi(x^2+\epsilon^2)}$
\item $F_n(x)=\dfrac{1-\cos{nx}}{\pi nx^2}$
\end{enumerate}

\begin{dimo}[Dimostriamo che la prima successione converge alla $\delta$ in $S'$]
\begin{align*}
F_n(\varphi)&=\dfrac{n}{\sqrt{\pi}}\int_{\mathbb{R}} e^{-n^2x^2}\varphi(x)\,dx \qquad \text{cambio di variabili }nx\rightarrow x \\
&= \dfrac{1}{\sqrt{\pi}} \int_{\mathbb{R}} e^{-x^2} \varphi\left(\frac{x}{n}\right)\, dx
\end{align*}
Si ha $\left|e^{-x^2}\varphi\left(\frac{x}{n}\right)\right|\leq e^{-x^2} ||\varphi||_{00}=g\in L^1$ quindi usando la convergenza dominata
\[
\lim_{n\rightarrow\infty}F_n(\varphi)=\dfrac{1}{\sqrt{\pi}} \int_{\mathbb{R}}\lim_{n\rightarrow\infty} e^{-x^2}\varphi\left(\frac{x}{n}\right)\,dx=\dfrac{\varphi(0)}{\sqrt{\pi}}\int_{\mathbb{R}} e^{-x^2}\,dx=\varphi(0) \qedhere
\] 
\end{dimo}


\subsection{Operazioni in $S'$}

Sia $A:S\rightarrow S$ $\varphi\mapsto A\varphi$ un'operazione continua e lineare su $S$. Vogliamo estenderla ad un'operazione su $S'$.

Supponiamo che esista un operatore lineare e continuo $A^*:S\rightarrow S$ tale che per ogni distribuzione regolare $\varphi_0\in S$ si abbia
\[
(A\varphi_0)(\varphi)=\varphi_0 (A^*\varphi)
\]
ovvero
\[
\int (A\varphi_0)(x)\,\varphi(x)\,d^Dx = \int \varphi_0(x)\,A^*\varphi(x)\,d^Dx
\]

Allora per ogni $F\in S'$ definiamo l'operatore $A:S'\rightarrow S'$ come 
\[
(AF)(\varphi)=F(A^*\varphi) \qquad \forall \varphi\in S
\]

\begin{itemize}[label=$\cdot$]
\item $AF$ è lineare e ben definito perchè $A^*\varphi\in S$
\item $AF\in S'$ perchè è definito come combinazione di mappe continue
\item La mappa $A:s'\rightarrow S'$ è continua perchè per $AF=\tilde{F}$
\[
|\tilde{F}(\varphi)|=|F(A^*\varphi)|=|F(\varphi^*)|<\epsilon
\]
\end{itemize}

Evidenziamo che in generale non è ben definito il prodotto $\psi\chi$ tra due distribuzioni, nel caso particolare in cui però una delle due distribuzioni $\psi=f$ sia regolare e anche in $O_M$ allora il prodotto $f\chi$ è ancora una distribuzione e 
\[
(f\chi)(\varphi)=\chi(f\varphi) \quad \forall \chi\in S'
\]


\subsubsection{Derivata di una distribuzione}

Ogni $F\in S'(\mathbb{R}^D$ ammette derivata parziale di ogni ordine in ogni direzione in senso distribuzionale e queste soddisfano $\partial_i\partial_j F=\partial_j\partial_i F$ $\forall F\in S',\forall i\in[1,D]$.

La derivata distribuzionale sarà una mappa $A:S'\rightarrow S'$ quindi la costruiamo usando una mappa continua in $S$ e definendo l'effetto di $A F$ su ogni funzione di test.

Usando una distribuzione regolare $F_0\in S$ avremo che la sua derivata $AF_0$ agisce su funzioni di test secondo
\[
(AF_0)(\varphi)=\int_{\mathbb{R}} \partial_j F_0(x)\, \varphi(x)\,dx = [F_0(x)\varphi(x)]_{-\infty}^{+\infty} -\int_{\mathbb{R}} F_0(x)\,\partial\varphi(x)\,dx=-\int_\mathbb{R} F_0(x)\,\partial\varphi(x)\,dx
\]
Dove nell'ultimo passaggio ricordiamo che le funzioni di test si annullano ad infinito estremamente velocemente. In $\mathbb{R}^D$ gli integrali sono $D-$dimensionali ma il risultato è lo stesso.

Questi passaggi si possono replicare formalmente senza nessun cambiamento per distribuzioni $F$ qualsiasi quindi definiamo la derivata distribuzionale come
\begin{equation}
(\partial_j F)(\varphi)=-F(\partial_j\varphi)
\end{equation}

Più in genereale se $P(\partial)$ è un polinomio nelle derivate si ha che
\[
\left(P(\partial)F\right)(\varphi)=F\left(P(-\partial)\varphi\right)
\]

\begin{esempio}[Esempi di derivate distribuzionali]
\begin{itemize}[label=$\cdot$]
\item La funzione di Heaviside ha come derivata la delta di Dirac.
\[
H'(x-a)(\varphi)=-H(x-a)\varphi'(x)=-\int_{-\infty}^{+\infty} H(x-a)\varphi'(x)\,dx=-\int_{a}^{+\infty} \varphi'(x)\,dx=\varphi(a)
\]
\item Le derivate della delta di Dirac sono:
\[
\left(\delta^{(N)}(x-a)\right)(\varphi)=\delta(x-a)\left((-)D^N \varphi\right)=-^N\delta(x-a)(D^N \varphi)
\]
quindi $\delta^{(N)}(x-a)=\delta(x-a)D^N$
\item $(\ln{|x|})'=\mathcal{P}(\sfrac{1}{x})$
\end{itemize}
\end{esempio}

In particolare la derivata distribuzionale di una distribuzione regolare rappresentata da un $f$ di classe $C^1$ coincide con la derivata ordinaria di $f$.

Dato che la derivata distribuzionale è una mappa continua in $S'$ commuterà con gli $S'-\lim$.

\subsubsection{Convoluzione}
g

\subsubsection{Composizione della $\delta$ con una funzione}
u

\section{Trasformata di Fourier}
\begin{definizione}
Definiamo la trasformata di Fourier $\mathcal{F}$ di un $\varphi\in L^1(\mathbb{R}^D)$ come
\begin{equation}
\hat{\varphi}(k)=
\dfrac{1}{(2\pi)^{\frac{D}{2}}}
\int_{\mathbb{R}^D}e^{-i\vec{k}\cdot\vec{x}}\varphi(x)\,d^Dx
\end{equation}
e l'antitrasformata di Fourier $\mathcal{F}^{-1}$ come
\begin{equation}
\check{\varphi}(k)=
\dfrac{1}{(2\pi)^{\frac{D}{2}}}
\int_{\mathbb{R}^D}e^{i\vec{k}\cdot\vec{x}}\varphi(x)\,d^Dx
\end{equation}
\end{definizione}

Ora vorremmo capire se per una $\varphi\in S$ la sua trasformata sia ancora in $S\subset L^1$, se la $\mathcal{F}:\varphi\mapsto\hat{\varphi}$ sia continua e invertibile e se in tal caso $\mathcal{F}^{-1}$ sia la sua inversa.


\begin{itemize}
\item Per $\varphi\in S(\mathbb{R}^D)$ la sua trasformata $\hat{\varphi}$ è di classe $C^\infty$
 \begin{dimo}[Lo mostraiamo  per $D=1$]
 \[
 \hat{\varphi}(k)=\dfrac{1}{\sqrt{2\pi}}\int e^{-ikx}\varphi(x)\,dx \quad \text{quindi}\quad
 \dfrac{ \hat{\varphi}(k+\Delta) -\hat{\varphi}(k)}{\Delta}=\dfrac{1}{\sqrt{2\pi}}\int e^{-ikx} \,  \dfrac{e^{-ix\Delta}-1}{\Delta x} \, x\varphi(x)\,dx
 \]
 E dato che $\left| e^{-ikx}\, \dfrac{e^{-ix\Delta}-1}{\Delta x} \, x\varphi(x)\right|\leq |x|\,|\varphi(x)|=g(x)\in L^1$ possiamo usare la convergenza dominata
 \begin{align*}
 \dfrac{\partial}{\partial k}\hat{\varphi}(k)&=\dfrac{1}{\sqrt{2\pi}}\int \lim_{\Delta\rightarrow0}  e^{-ikx}\, \dfrac{e^{-ix\Delta}-1}{\Delta x} \, x\varphi(x) = \dfrac{1}{\sqrt{2\pi}}\int e^{-ikx}-ix\varphi(x)\,dx=\\
 &=-i\mathcal{F}\left(x\varphi(x)\right)\,(k) \implies \hat{\varphi}\in C^\infty \qedhere
 \end{align*}
 \end{dimo}
\item Dato $P(x^m,\partial_i)$ un generico polinomio in $\vec{x}$ e $\vec{\nabla}$ vale
\begin{equation}
\mathcal{F}\left( P(x^m,\partial_i)\varphi \right) \, (\vec{k}) =
P\left(i\dfrac{\partial}{\partial k^m},ik^i\right)\hat{\varphi}(\vec{k})
\end{equation}
\item Per  $\varphi\in S$ si ha che $\hat{\varphi}\in S$ e la trasformata di Fourier è una mappa continua da $S$ in $S$
\begin{dimo}
\begin{align*}
\left| k^N\partial^M \hat{\varphi}(k) \right| &=\left| \dfrac{\partial^N}{\partial x^N} \left( x^M \varphi(x)\right) \, (k)\right| = \dfrac{1}{\sqrt{2\pi}}\left| \int e^{-ikx} \partial^N\left(x^M\varphi(x)\right)\,dx \right| \\ &\leq \dfrac{1}{\sqrt{2\pi}} \int \left| \partial^N\left(x^M\varphi(x)\right)\right|\,dx\\
\left|\left|\hat{\varphi}\right|\right|_{NM}&=\sup_{k\in\mathbb{R}}\left| k^N\partial^M\hat{\varphi}(k) \right|\leq\dfrac{1}{\sqrt{2\pi}}\int \left| \partial^N\left(x^M\varphi(x)\right)\right|\,dx\\
&\leq \dfrac{1}{\sqrt{2\pi}}\int \dfrac{1+x^2}{1+x^2}\left| \partial^N\left(x^M\varphi(x)\right)\right|\,dx\\&\leq C\sum_{\text{finita}}||\varphi|| \qquad \text{non entriamo troppo nei dettagli}\qedhere
\end{align*}
\end{dimo}
\end{itemize} 

\begin{proposizione}$\mathcal{F}^{-1}\circ\mathcal{F}=\mathcal{F}\circ\mathcal{F}^{-1}=\mathbb{I}_{S}$\end{proposizione}
\begin{dimo}

\end{dimo}

Le proprietà della trasformata di Fourier in $S$ sono:\begin{itemize}[label=$\cdot$]
\item $\hat{\hat{\varphi}}(x)=\varphi(-x)$ 
\item $\mathcal{F}^4=\mathbb{I}$
\item $\mathcal{F}\left(\varphi(x+a)\right)\,(k)=e^{ika}\hat{\varphi}(k)$
\item $\mathcal{F}\left(e^{i\alpha x}\varphi(x)\right)\,(k)=\hat{\varphi}(k-a)$
\item $\mathcal{F}\left(\varphi_1 * \varphi_2\right) (k)=(2\pi)^{\frac{D}{2}}\hat{\varphi_1}(k)\hat{\varphi_2}(k)$
\item $\mathcal{F}(\varphi_1 \varphi_2)\,(k)=\dfrac{1}{(2\pi)^\frac{D}{2}}(\hat{\varphi_1} * \hat{\varphi_2})(k)$
\end{itemize}

 
 
\subsection{Tasformata di Fourier in $S'$}
`

\subsection{Trasformata di Fourier ed equazioni differenziali alle derivate parziali}
;

\section{Operatori su spazi di Hilbert}
l

\subsection{Derivata}





































































































































































































































































































































\end{document}